% ==============================================================================
% Chapter 10: THE CALCULUS OF CUSPS AND COMPONENTS VIA THE GROUPS T[N],
%             AND THE GLOBAL STRUCTURE OF THE BASIC MODULI PROBLEMS
% ==============================================================================
\chapter{The Calculus of Cusps and Components via the Groups $\mathbf{T}[N]$, and the Global Structure of the Basic Moduli Problems}
\label{ch:10}

% Section 10.1
\section{Motivation}
\label{sec:10.1}

We have already seen that the problem of analyzing the structure near infinity of a relatively representable $\mathcal{P}$ on $\Ell/R$ is essentially the problem of ``computing'' the $R((q))$-scheme
\[
\mathcal{P}_{\Tate(q)/R((q))} ,
\]
with its action of the group $\pm 1 = \Aut(\Tate(q)/R((q)))$.
When $\mathcal{P}$ is a problem of level $N$, i.e., when $\mathcal{P} = \mathcal{F} \circ [N]$ for some relatively representable moduli problem $\mathcal{F}$ on $(\FLFG/R)$, (cf.~\ref{item:4.14.2}) then we have
\[
\mathcal{P}_{\Tate(q)/R((q))} = \mathcal{F}_{\Tate(q)[N]/R((q))}
\]
with $\pm 1$ operating on $\mathcal{F}$ as the automorphism $\pm 1$ of $\Tate(q)[N]/\allowbreak R((q))$.
We have already seen that for any integer $N \ge 1$, the group $\Tate(q)[N]$ over $R((q))$ with its $e_N$-pairing, is isomorphic to the group
\[
\mathbf{T}[N] \tens{\Z[q,q^{-1}]} R((q)) , \quad \text{with its } e_N\text{-pairing}
\]
obtained from $\mathbf{T}[N]/\Z[q,q^{-1}]$ by extension of scalars. Thus we have
\[
\mathcal{P}_{\Tate(q)/R((q))} = (\mathcal{F}_{\mathbf{T}[N]/R[q,q^{-1}]}) \tens{R[q,q^{-1}]} R((q)) .
\]
In this chapter we will explicitly calculate
\[
\mathcal{F}_{\mathbf{T}[N]/R[q,q^{-1}]}
\]
for various mixes of the basic moduli problems over $\Z$ and their ``canonical'' incarnations as moduli problems over cyclotomic integer rings. As explained above, these calculations will reveal the structure near infinity of the moduli problems in question (cf.~\ref{sec:10.8}).

By combining this information with two other ingredients,
\begin{enumerate}[label=(\arabic*)]
\item the homogeneity principle, (which already played so large a role in Chapter~5),
\item the \emph{transcendental} description of our moduli spaces as quotients of the upper-half plane, (a description which we have up to now avoided), used in the proof of \ref{cor:10.9.2},
\end{enumerate}
we obtain a good picture of the global structure of our moduli schemes (cf.~\ref{sec:10.9} through \ref{sec:10.13}).

We will see later, in Chapter~13, that the calculations (\ref{sec:10.2}--\ref{sec:10.8}) of this chapter also reveal the structure of the reductions $\bmod p$ of these same moduli problems.


% Section 10.2
\section{Analysis of $[\Gamma(N)]$}
\label{sec:10.2}

\kmitem\label{item:10.2.1}
Over any base $S$, consider a finite locally free commutative group-scheme $G/S$ of rank $N^2$ which is killed by $N$, and which is given as an extension
\[
0 \longrightarrow \mu_N \longrightarrow G \xrightarrow{\ \pi\ } \Z/N\Z \longrightarrow 0 .
\]
Suppose we are given a $\Gamma(N)$-structure on $G/S$, which we will view as a homomorphism of $S$-groups
\[
\phi : (\Z/N\Z)^2 \to G
\]
which is a ``$(\Z/N\Z)^2$-generator'' of $G$ (cf.~\ref{km:1.10.5}). The corresponding Drinfeld $N$-basis $(P,Q)$ of $G/S$ is related to $\phi$ by
\[
\phi\begin{pmatrix} 1 \\ 0 \end{pmatrix} = P , \quad \phi\begin{pmatrix} 0 \\ 1 \end{pmatrix} = Q
\]
\[
\phi\begin{pmatrix} x \\ y \end{pmatrix} = x P + y Q = (P,Q)\begin{pmatrix} x \\ y \end{pmatrix} .
\]
If more than one $\Gamma(N)$-structure is in play, we will write
\[
(P_\phi, Q_\phi)
\]
for the Drinfeld $N$-basis corresponding to a particular $\phi$.

\begin{klem}\label{lem:10.2.2}
The composite homomorphism $\Lambda = \pi \cdot \phi$
\begin{equation*}
\begin{tikzcd}[row sep=2em, column sep=2em]
& (\Z/N\Z)^2 \arrow[ld, "\phi"'] \arrow[rd, "\Lambda = \pi \cdot \phi"] & \\
0 \arrow[r] & \mu_N \arrow[r] & G \arrow[r, "\pi"] & \Z/N\Z \arrow[r] & 0
\end{tikzcd}
\end{equation*}
is surjective, and its kernel $\Ker(\Lambda) \subset (\Z/N\Z)^2$ is canonically isomorphic to $\Z/N\Z$, with canonical generator $k_\Lambda$.
\end{klem}

\begin{proof}
That $\Lambda$ is surjective is a special case of \eqref{prop:1.11.2}. Its kernel $\Ker(\Lambda)$ is therefore a finite etale subgroup of $(\Z/N\Z)^2$, so constant on each connected component of $S$. The quotient group $(\Z/N\Z)^2/\Ker(\Lambda)$ is given as isomorphic to $\Z/N\Z$. Therefore $\Ker(\Lambda)$ must be a cyclic group of order $N$ (elementary divisors), and it has a canonical generator $k_\Lambda$ determined by the requirement that for any vector $\ell \in (\Z/N\Z)^2$,
\[
\det(k_\Lambda, \ell) = \Lambda(\ell) . \qedhere
\]
\end{proof}

\kmitem\label{item:10.2.3}
In view of the lemma, we can uniquely complete the above diagram to a commutative diagram
\begin{equation}\label{eq:10.2.3.1}
\begin{tikzcd}[row sep=2em, column sep=2em]
0 \arrow[r] & \Z/N\Z \arrow[r, "k_\Lambda"] \arrow[d, "\phi(k_\Lambda)"'] & (\Z/N\Z)^2 \arrow[r, "\Lambda"] \arrow[d, "\phi"'] & \Z/N\Z \arrow[r] \arrow[d, equal] & 0 \\
0 \arrow[r] & \mu_N \arrow[r] & G \arrow[r, "\pi"] & \Z/N\Z \arrow[r] & 0 .
\end{tikzcd}
\end{equation}
According to \eqref{prop:1.11.2}, $\phi(k_\Lambda)$ will necessarily be a generator of $\mu_N$.

\begin{kprop}\label{prop:10.2.4}
For each $\Lambda$ in $\HomSurj((\Z/N\Z)^2, \Z/N\Z)$, let $k_\Lambda$ denote the canonical generator of $\Ker(\Lambda)$, defined by the requirement that
\[
\det(k_\Lambda, \ell) = \Lambda(\ell)
\]
for all $\ell \in (\Z/N\Z)^2$. Choose a vector
\[
\ell_\Lambda \in (\Z/N\Z)^2 ,
\]
such that
\[
\Lambda(\ell_\Lambda) = 1 .
\]
Let $G/S$ be as above, with $S$ connected. Then
\begin{enumerate}[label=(\arabic*)]
\item If $\phi$ is a $\Gamma(N)$-structure on $G/S$, then $\pi \cdot \phi = \Lambda$ lies in $\HomSurj((\Z/N\Z)^2, \Z/N\Z)$, $\phi(k_\Lambda)$ is a primitive $N$'th root of unity in $\mu_N(S)$, and $\phi(\ell_\Lambda) \in G(S)$ has $\pi\phi(\ell_\Lambda) = 1 \in \Z/N\Z$.
\item Given $\Lambda \in \HomSurj((\Z/N\Z)^2, \Z/N\Z)$, consider the oriented basis $k_\Lambda$, $\ell_\Lambda$ of $(\Z/N\Z)^2$. A homomorphism $\phi$ of $S$-groups
\begin{equation*}
\begin{tikzcd}[row sep=1.8em]
(\Z/N\Z)^2 \arrow[d, "\phi"'] \\
0 \arrow[r] & \mu_N \arrow[r] & G \arrow[r, "\pi"] & \Z/N\Z \arrow[r] & 0
\end{tikzcd}
\end{equation*}
is a $\Gamma(N)$-structure with $\pi \cdot \phi = \Lambda$ if and only if both of the following conditions are satisfied:
\[
\resizebox{0.96\linewidth}{!}{$\displaystyle
\left\{\begin{aligned}
& \phi(k_\Lambda) \text{ lies in } \mu_N(S), \text{ and is a primitive } N\text{'th root of unity.} \\
& \phi(\ell_\Lambda) \text{ lies in } G(S), \text{ and } \pi(\phi(\ell_\Lambda)) = 1 \text{ in } \Z/N\Z.
\end{aligned}\right.
$}
\]
\end{enumerate}
\end{kprop}

\begin{proof}
Immediate from \eqref{prop:1.11.2}.
\end{proof}

\begin{kcor}\label{cor:10.2.5}
Let $S$ be an arbitrary scheme, $G/S$ a finite locally free $S$-group killed by $N$ given with an f.p.p.f.\ exact sequence
\[
0 \longrightarrow \mu_N \longrightarrow G \xrightarrow{\ \pi\ } \Z/N\Z \longrightarrow 0 .
\]
Let $G_1 \subset G$ denote the closed subscheme $\pi^{-1}(1)$, so that $G_1$ is a $\mu_N$-torsor on $S$. Then we have a canonical isomorphism of $S$-schemes
\[
[\Gamma(N)]_{G/S} \xrightarrow{\ \sim\ } \coprod_{\Lambda \in \HomSurj((\Z/N\Z)^2, \Z/N\Z)} \Z[\zeta_N] \tens{\Z} G_1 ,
\]
whose description in terms of $T$-valued points, for any connected $S$-scheme $T$, is given by
\[
\phi \mapsto (\pi \cdot \phi, \phi(k_\Lambda), \phi(\ell_\Lambda)) .
\]
The $\Gamma(N)$-structure
\[
(\Lambda, \zeta, X \in G_1(T)) ,
\]
with
\[
\left\{\begin{aligned}
\phi(k_\Lambda) & = \zeta \\
\phi(\ell_\Lambda) & = X
\end{aligned}\right.
\]
will sometimes be denoted
\[
\phi \leftrightarrow (\Lambda_\phi, \zeta_\phi, X_\phi) .
\]
\end{kcor}

% Section 10.3
\section{Group action}
\label{sec:10.3}

We now give the explicit description of the right action of the group $\GL(2,\Z/N\Z) = \Aut((\Z/N\Z)^2)$ on the scheme $[\Gamma(N)]_{G/S}$. The action is defined by
\begin{equation}\label{eq:10.3.1}
(\phi, g) \mapsto \phi \circ g .
\end{equation}

% The display (10.3.1) sits directly under the section (no numbered item precedes
% it) and occupies the first slot of the per-section sequence; we advance the
% kmitem counter by hand so that Lemma 10.3.2 continues the same sequence.
\addtocounter{kmitem}{1}

\begin{klem}\label{lem:10.3.2}
Under the right action $\phi \mapsto \phi \circ g$ of $\GL(2,\Z/N\Z)$ on $[\Gamma(N)]_{G/S}$, we have
\begin{enumerate}[label=(\arabic*)]
\item $\Lambda_{\phi g} = \Lambda_\phi \circ g$,
\item $\ell_{\Lambda \cdot g} = g^{-1}(\ell_\Lambda) + n(g,\Lambda) g^{-1}(k_\Lambda)$ for some $n(g,\Lambda) \in \Z/N\Z$,
\item $k_{\Lambda \cdot g} = \det(g) g^{-1}(k_\Lambda)$, $\Ker(\Lambda g) = g^{-1}(\Ker(\Lambda))$,
\item $\zeta_{\phi g} = (\zeta_\phi)^{\det(g)}$,
\item $X_{\phi g} = (\zeta_\phi)^{n(g,\Lambda)} X_\phi$.
\end{enumerate}
\end{klem}

\begin{proof}
By definition, $\Lambda_\phi = \pi \circ \phi$, whence $\Lambda_{\phi g} = \pi \circ \phi \circ g = \Lambda_\phi \circ g$. By definition, $\ell_\Lambda$ is a chosen vector one for each $\Lambda$, with $\Lambda(\ell_\Lambda) = 1$. Therefore $g^{-1}(\ell_\Lambda)$ is a candidate for $\ell_{\Lambda \cdot g}$, so their difference must lie in $\Ker(\Lambda g) = g^{-1}\Ker(\Lambda)$. The vector $k_{\Lambda g}$ is characterized by the properties
\[
\det(k_{\Lambda g}, \ell_{\Lambda g}) = 1, \quad k_{\Lambda g} \in \Ker(\Lambda g) .
\]
Therefore $k_{\Lambda g}$ is a $(\Z/N\Z)^\times$-multiple of the obvious generator $g^{-1}(k_\Lambda)$ of $\Ker(\Lambda g)$, say $k_{\Lambda g} = a g^{-1}(k_\Lambda)$. To compute $a$, we write
\begin{align*}
1 = \det(k_{\Lambda g}, \ell_{\Lambda g}) & = \det(a g^{-1}(k_\Lambda), g^{-1}(\ell_\Lambda)) \\
& = a \det(g)^{-1} \det(k_\Lambda, \ell_\Lambda) \\
& = a \det(g)^{-1} .
\end{align*}
By definition
\begin{align*}
\zeta_{\phi g} = \phi g(k_{\Lambda g}) & = \phi g(\det(g) g^{-1}(k_\Lambda)) \\
& = \phi(k_\Lambda)^{\det(g)} \\
& = (\zeta_\phi)^{\det(g)} .
\end{align*}
Similarly, we have
\begin{align*}
X_{\phi g} = \phi g(\ell_{\Lambda \cdot g}) & = \phi g(g^{-1}(\ell_\Lambda) + n(g,\Lambda) g^{-1}(k_\Lambda)) \\
& = \phi(\ell_\Lambda + n(g,\Lambda) k_\Lambda) \\
& = (\zeta_\phi)^{n(g,\Lambda)} X_\phi . \qedhere
\end{align*}
\end{proof}

\kmitem\label{item:10.3.3}
For each $\Lambda \in \HomSurj((\Z/N\Z)^2, \Z/N\Z)$, let $\Fix(\Lambda)$ be the subgroup of $\GL(2,\Z/N\Z)$ consisting of all $g$ such that $\Lambda g = \Lambda$. This group acts stably on the $\Lambda$-component of $[\Gamma(N)]_{G/S}$,
\[
\Z[\zeta_N] \tens{\Z} G_1 ,
\]
by
\[
(\zeta, X) \mapsto (\zeta^{\det(g)}, \zeta^{n(g,\Lambda)} X) .
\]

\kmitem\label{item:10.3.4}
The group $\Fix(\Lambda)$ is isomorphic to the subgroup $\begin{pmatrix} * & * \\ 0 & 1 \end{pmatrix}$ of $\GL(2,\Z/N\Z)$, by the construction
\[
g \in \Fix(\Lambda) \mapsto \begin{pmatrix} \det(g) & n(g,\Lambda) \\ 0 & 1 \end{pmatrix} .
\]
[Of course, for $\Lambda$ the homomorphism $\begin{pmatrix} x \\ y \end{pmatrix} \mapsto y = (0,1)\begin{pmatrix} x \\ y \end{pmatrix}$, $\Fix(\Lambda)$ is \emph{already} the subgroup $\begin{pmatrix} * & * \\ 0 & 1 \end{pmatrix}$ of $\GL(2,\Z/N\Z)$; our $(k_\Lambda, \ell_\Lambda)$-construction is just an explicit incarnation of the fact that $\GL(2,\Z/N\Z)$ acts transitively on $\HomSurj$, so that every $\Fix(\Lambda)$ is conjugate to the ``standard'' one. From this point of view, the set $\HomSurj((\Z/N\Z)^2, \Z/N\Z)$ is the homogeneous space
\begin{equation}\label{eq:10.3.4.1}
\begin{pmatrix} * & * \\ 0 & 1 \end{pmatrix} \Big\backslash \GL(2,\Z/N\Z) \xrightarrow{\ \sim\ } \HomSurj((\Z/N\Z)^2, \Z/N\Z)
\end{equation}
if we take $\Lambda = (0,1)$ as base point.]

\begin{kcor}\label{cor:10.3.5}
Let $\Gamma$ be an arbitrary subgroup of $\GL(2,\allowbreak \Z/N\Z)$. Let $\{\Lambda_i\}$ be a set of representatives in $\HomSurj((\Z/N\Z)^2,\allowbreak \Z/N\Z)$ of the orbit space
\[
\HomSurj/\Gamma
\]
for the right action $(\Lambda \mapsto \Lambda \circ g)$ of $\Gamma$ on $\HomSurj((\Z/N\Z)^2, \Z/N\Z)$.

Let $S$ be an arbitrary scheme, $G/S$ a finite locally free $S$-group killed by $N$ given with an f.p.p.f.\ exact sequence
\[
0 \longrightarrow \mu_N \longrightarrow G \xrightarrow{\ \pi\ } \Z/N\Z \longrightarrow 0 .
\]
Then the quotient of the $S$-scheme $[\Gamma(N)]_{G/S}$ by the given subgroup $\Gamma$ of $\GL(2,\Z/N\Z)$ is given by a canonical isomorphism
\[
([\Gamma(N)]_{G/S})/\Gamma \simeq \coprod_{\Lambda_i \in \HomSurj/\Gamma} (\Z[\zeta_N] \tens{\Z} G_1)/(\Gamma \cap \Fix(\Lambda_i)) .
\]
\end{kcor}


% Section 10.4
\section{Canonical problems}
\label{sec:10.4}

\begin{klem}\label{lem:10.4.1}
Let $G/S$ be a finite locally free commutative group-scheme of rank $N^2$, killed by $N$, given with a strongly alternating autoduality
\[
(\ ,\ ) : G \times G \to \mu_N ,
\]
which after some faithfully flat base change $S' \to S$ becomes isomorphic to $(E'[N], e_N)$ for some elliptic curve $E'/S'$. Then for any $\Gamma(N)$-structure $\phi = (P,Q)$ on $G/S$, the pairing-value
\[
(P,Q) \in \mu_N(S)
\]
is a primitive $N$'th root of unity on $S$.
\end{klem}

\begin{proof}
Because $S'$ is faithfully flat over $S$, to verify that $(P,Q)$ is a zero of $\Phi_N(X)$ in $\Gamma(S, \cO_S)$ it suffices to do so over $S'$, where the result is already known \eqref{thm:5.6.3}.
\end{proof}

\kmitem\label{item:10.4.2}
Let $H$ be a subgroup of $(\Z/N\Z)^\times$ and $A$ the corresponding subring of $\Z[\zeta_N]$:
\[
A = (\Z[\zeta_N])^H \subset \Z[\zeta_N] .
\]
Let $S$ be any $A$-scheme, $(G/S, (\ ,\ ))$ as in the above lemma. We say that a $\Gamma(N)$-structure $\phi = (P,Q)$ on $G/S$ is ``$A$-canonical'' if the structure of $\Z[\zeta_N]$-scheme which $S$ acquires from the primitive $N$'th root of unity
\[
(P,Q) \in \mu_N^\times(S)
\]
is compatible with its given structure of $A$-scheme, and with the given inclusion of $A$ in $\Z[\zeta_N]$; in other words, $A$-canonicity of $(P,Q)$ means that the diagram
\begin{equation*}
\begin{tikzcd}[row sep=2em, column sep=1.5em]
& S \arrow[ld, "\text{given}"'] \arrow[rd, "{\text{via } (P,Q)}"] & \\
\Spec(A) & & \Spec(\Z[\zeta_N]) \arrow[ll, "\text{given}"]
\end{tikzcd}
\end{equation*}
commutes.

We denote by $[\Gamma(N)]^{A\text{-can}}$ the corresponding moduli problem on $\Ell/A$.

\begin{kthm}\label{thm:10.4.3}
Let $H$ be a subgroup of $(\Z/N\Z)^\times$, $A$ the subring $(\Z[\zeta_N])^H$ of $\Z[\zeta_N]$. Let $S$ be an arbitrary $A$-scheme, and $G/S$ a finite locally free commutative $S$-group of rank $N^2$, killed by $N$, given with an f.p.p.f.\ short exact sequence
\[
0 \longrightarrow \mu_N \longrightarrow G \longrightarrow \Z/N\Z \longrightarrow 0 .
\]
Recall that $G/S$ carries a unique strongly alternating autoduality $(\ ,\ )$ whose associated $\langle\ ,\ \rangle : \mu_N \times \Z/N\Z \to \mu_N$ is the standard pairing $\langle \zeta, a \rangle = \zeta^a$ (cf.~\ref{lem:8.10.4}), so we can speak of $A$-canonical $\Gamma(N)$-structures on $(G/S, (\ ,\ ))$.
\begin{enumerate}[label=(\arabic*)]
\item We have a canonical isomorphism of $S$-schemes
\[
([\Gamma(N)]^{A\text{-can}})_{G/S} = \coprod_{\HomSurj((\Z/N\Z)^2, \Z/N\Z)} \Z[\zeta_N] \tens{A} G_1 .
\]
\item For any subgroup $\Gamma \subset \GL(2,\Z/N\Z)$ such that
\[
\det(\Gamma) \subset H ,
\]
let $\{\Lambda_i\}$ be a set of representatives for the orbit space
\[
\HomSurj/\Gamma .
\]
Then we have a canonical isomorphism of $S$-schemes
\[
\resizebox{0.97\linewidth}{!}{$\displaystyle
(([\Gamma(N)]^{A\text{-can}})_{G/S})/\Gamma \simeq \coprod_{\Lambda_i \in \HomSurj/\Gamma} (\Z[\zeta_N] \tens{A} G_1)/(\Gamma \cap \Fix(\Lambda_i)) .
$}
\]
\end{enumerate}
\end{kthm}

\begin{proof}
Exactly the same as for its $\Z$-analogues \eqref{cor:10.2.5} and \eqref{cor:10.3.5}.
\end{proof}


% Section 10.5
\section{Explication for $\mathbf{T}[N]$}
\label{sec:10.5}

\begin{kthm}\label{thm:10.5.1}
Let $H$ be a subgroup of $(\Z/N\Z)^\times$, $A$ the subring $\Z[\zeta_N]^H$ of $\Z[\zeta_N]$. For any $A[q,q^{-1}]$-algebra $B$, denote by $\mathbf{T}[N]/B$ the group-scheme
\[
\mathbf{T}[N] \tens{\Z[q,q^{-1}]} B ,
\]
with its canonical extension structure and $e_N$-pairing. Then
\begin{enumerate}[label=(\arabic*)]
\item we have a canonical isomorphism of $B$-schemes
\[
\resizebox{0.96\linewidth}{!}{$\displaystyle
([\Gamma(N)]^{A\text{-can}})_{\mathbf{T}[N]/B} = \coprod_{\HomSurj((\Z/N\Z)^2, \Z/N\Z)} \Spec\bigl(\Z[\zeta_N] \tens{A} (B[X]/(X^N = q))\bigr) .
$}
\]
\item For any subgroup $\Gamma \subset \GL(2,\Z/N\Z)$ with $\det(\Gamma) \subset H$, let $\{\Lambda_i\}$ be a set of representatives for $\HomSurj/\Gamma$. Then we have a canonical isomorphism of $B$-schemes
\[
\resizebox{0.96\linewidth}{!}{$\displaystyle
(([\Gamma(N)]^{A\text{-can}})_{\mathbf{T}[N]/B})/\Gamma = \coprod_{\Lambda_i \in \HomSurj((\Z/N\Z)^2, \Z/N\Z)/\Gamma} \Spec\bigl((\Z[\zeta_N] \tens{A} (B[X]/(X^N = q)))^{\Gamma \cap \Fix(\Lambda_i)}\bigr) .
$}
\]
\end{enumerate}
\end{kthm}

\begin{proof}
For $G = \mathbf{T}[N]$, the $\mu_N$-torsor $G_1$, cf.~\eqref{cor:10.2.5}, is just $X^N = q$, and we apply \eqref{thm:10.4.3}.
\end{proof}


% Section 10.6
\section{Cusp-labels and component-labels}
\label{sec:10.6}

Let $\Gamma$ be a subgroup of $\GL(2,\Z/N\Z)$. The orbit space
\[
\HomSurj((\Z/N\Z)^2, \Z/N\Z)/\Gamma
\]
we propose to call the space of cusp-labels for $[\Gamma(N)]/\Gamma$, or, \emph{sometimes}, for $\Gamma$. If we fix the base-point $\Lambda\begin{pmatrix} x \\ y \end{pmatrix} = y = (0,1)\begin{pmatrix} x \\ y \end{pmatrix}$ for $\HomSurj$, this space of cusp-labels becomes the double coset space
\[
\begin{pmatrix} * & * \\ 0 & 1 \end{pmatrix} \Big\backslash \GL(2,\Z/N\Z)/\Gamma .
\]

If we view $(\Z/N\Z)^\times$ as the central subgroup of \emph{scalars} of $\GL(2,\Z/N\Z)$, then we can form the orbit space
\[
(\Z/N\Z)^\times \Big\backslash \HomSurj((\Z/N\Z)^2, \Z/N\Z)/\Gamma .
\]
As a double coset space, via base-point $\Lambda = (0,1)$, this space is
\[
\begin{pmatrix} * & * \\ 0 & * \end{pmatrix} \Big\backslash \GL(2,\Z/N\Z)/\Gamma .
\]
We call this orbit space the space of \emph{component-labels} for $\Gamma$, or for $[\Gamma(N)]/\Gamma$. There is a natural projection
\begin{equation*}
\begin{tikzcd}[row sep=4em]
\text{cusp-labels} \arrow[d] \\
\text{component-labels} ,
\end{tikzcd}
\end{equation*}
whose geometrical significance will become clear later when we reduce $\bmod\ p$, and which will explain the ``component-label'' terminology.

% Section 10.7
\section{Some combinatorial lemmas}
\label{sec:10.7}

\begin{klem}\label{lem:10.7.1}
If $N \ge 3$, then $-1$ acts without fixed points on $\HomSurj((\Z/N\Z)^2, \Z/N\Z)$.
\end{klem}

\begin{proof}
If $\Lambda = (a,b)$ is fixed by $-1$, $2a = 2b = 0$. As $(a,b)$ is primitive, $2 = 0$ in $\Z/N\Z$.
\end{proof}

\begin{klem}\label{lem:10.7.2}
Let $N$ be an odd integer, and $\Gamma$ a subgroup of $\SL(2,\Z/N\Z)$. Then either $\Gamma$ contains the scalar $-1$, or $-1$ operates without fixed points on the space $\HomSurj/\Gamma$ of cusp-labels for $\Gamma$.
\end{klem}

\begin{proof}
Suppose that $-1$ fixes the coset $\Lambda\Gamma$ for some $\Lambda \in \HomSurj$. This means that
\[
-\Lambda = \Lambda \circ \gamma \quad\text{for some}\quad \gamma \in \Gamma .
\]
In the basis $k_\Lambda$, $\ell_\Lambda$ of $(\Z/N\Z)^2$, the matrix of $\gamma$ is of the form
\[
\begin{pmatrix} * & * \\ 0 & -1 \end{pmatrix} .
\]
Because $\gamma \in \Gamma \subset \SL$, the matrix must be of the form
\[
\gamma = \begin{pmatrix} -1 & -x \\ 0 & -1 \end{pmatrix}
\]
for some $x \in \Z/N\Z$. Then
\[
\gamma^2 = \begin{pmatrix} 1 & 2x \\ 0 & 1 \end{pmatrix} .
\]
Because $N$ is odd, there exists $n \ge 1$ with $2n \equiv 1 \bmod N$, ($n = (N+1)/2$ for example). So
\[
\gamma^{2n} = \begin{pmatrix} 1 & 2nx \\ 0 & 1 \end{pmatrix} = \begin{pmatrix} 1 & x \\ 0 & 1 \end{pmatrix} = -\gamma ,
\]
and therefore $-1$ lies in $\Gamma$.
\end{proof}

\begin{klem}\label{lem:10.7.3}
Let $N = N_1 N_2$ with $N_1$ and $N_2$ relatively prime integers. Let $\Gamma$ be a subgroup of $\GL(2,\Z/N\Z)$ which is of the form $\Gamma_1 \times \Gamma_2$, with $\Gamma_i$ a subgroup of $\GL(2,\Z/N_i\Z)$, $i = 1,2$, when we view $\GL(2,\Z/N\Z)$ as the product of the $\GL(2,\Z/N_i\Z)$. Then:
\begin{enumerate}[label=(\arabic*)]
\item We have a natural bijection
\[
\resizebox{0.95\linewidth}{!}{$\displaystyle
\text{(cusp-labels for } \Gamma) \xrightarrow{\ \sim\ } \text{(cusp-labels for } \Gamma_1) \times \text{(cusp-labels for } \Gamma_2)
$}
\]
under which any scalar $a \in (\Z/N\Z)^\times$ acts (by right multiplication) as $(a \bmod N_1, \ a \bmod N_2)$.
\item If $-1$ operates without fixed points on the cusp-labels of $\Gamma_1$, then it operates without fixed points on the cusp-labels of $\Gamma$.
\item If $-1$ operates without fixed points on the cusp-labels of $\Gamma$, then for some $i \in \{1,2\}$, it operates without fixed points on the cusp-labels for $\Gamma_i$.
\end{enumerate}
\end{klem}

\begin{proof}
(1) is the Chinese Remainder Theorem, (2),(3) the observation that the only fixed points of a product endomorphism of a product space are products of fixed points.
\end{proof}

\begin{klem}\label{lem:10.7.4}
If $N \ne 1,2,4$, then for any subgroup $\Gamma$ of the standard upper unipotent subgroup
\[
\Gamma \subset \begin{pmatrix} 1 & * \\ 0 & 1 \end{pmatrix} \subset \GL(2,\Z/N\Z) ,
\]
the element $-1$ acts without fixed points on the space of cusp-labels of $\Gamma$.
\end{klem}

\begin{proof}
Visibly $-1 \notin \Gamma$. So for odd $N$, this is a particular case of \eqref{lem:10.7.2}.

In the remaining cases, we argue as follows. The existence of a fixed point of $-1$ acting on $\HomSurj/\Gamma$ means that there exists some $\Lambda \in \HomSurj$, and some $\gamma \in \Gamma$, with
\[
-\Lambda = \Lambda \gamma .
\]
Matricially, this is
\[
-(x,y) = (x,y)\begin{pmatrix} 1 & B \\ 0 & 1 \end{pmatrix} = (x, Bx+y)
\]
for some primitive vector $(x,y) \in (\Z/N\Z)^2$, and some $B \in \Z/N\Z$. Thus we have
\[
2x \equiv 0 , \quad 2y \equiv -Bx \quad \bmod N .
\]
If $N = 2M$ with $M$ odd and $M \ge 3$, then we infer
\[
x \equiv 0(M), \quad 2y \equiv -Bx \equiv 0(M) .
\]
Because $M$ is odd, we may infer $y \equiv 0(M)$ also. Thus $(x,y)$ is not primitive, contradiction.

If $N = 4M$ with $M \ge 2$ arbitrary, we again consider
\[
2x \equiv 0(4M), \quad 2y \equiv -Bx \ (4M) .
\]
The first congruence shows $x \equiv 0(2M)$, $4y \equiv -2Bx \equiv 0(4M)$ so $y \equiv 0(M)$. But $x \equiv 0(2M)$ implies $x$ is nilpotent mod $4M$. As $(x,y)$ is primitive mod $4M$, this implies that $y$ is a unit mod $4M$, hence that $y \bmod M$ is a unit mod $M$, contradicting $y \equiv 0(M)$.
\end{proof}


% Section 10.8
\section{Application to structure near infinity}
\label{sec:10.8}

\kmitem\label{item:10.8.1}
Let $H$ be a subgroup of $(\Z/N\Z)^\times$, $A$ the subring $\Z[\zeta_N]^H$ of $\Z[\zeta_N]$. We know \eqref{thm:5.1.1} that $[\Gamma(N)]$ is a regular moduli problem, finite and flat over $\Ell/\Z$. By (\ref{cor:9.1.8}, \ref{cor:9.1.9}), we know that $[\Gamma(N)]^{A\text{-can}}$ is a regular moduli problem, finite and flat over $\Ell/A$. For any subgroup $\Gamma \subset \GL(2,\Z/N\Z)$ with $\det(\Gamma) \subset H$, we know that $\Gamma$ acts on $[\Gamma(N)]^{A\text{-can}}$ as moduli problem on $\Ell/A$, and that the quotient $[\Gamma(N)]^{A\text{-can}}/\Gamma$ is itself a normal moduli problem, finite over $\Ell/A$, and canonically isomorphic to $([\Gamma(N)]/\Gamma)^{A\text{-can}}$.

\begin{kthm}\label{thm:10.8.2}
Let $H$ be a subgroup of $(\Z/N\Z)^\times$, $A$ the subring $\Z[\zeta_N]^H$ of $\Z[\zeta_N]$, $B$ a noetherian $A$-algebra. Let $\Gamma \subset \GL(2,\Z/N\Z)$ be a subgroup with $\det(\Gamma) \subset H$, and denote by $\mathcal{P}$ the moduli problem
\[
\mathcal{P} = ([\Gamma(N)]/\Gamma)^{A\text{-can}} \tens{A} B \quad\text{on } \Ell/B \; .
\]
Then we have a canonical isomorphism of $B((q))$-schemes
\begin{gather*}
\mathcal{P}_{\Tate(q)/B((q))} \simeq \\
\resizebox{\linewidth}{!}{$\displaystyle
\coprod_{\substack{\Lambda_i\ \text{reps of}\\ \HomSurj/\Gamma}}
\Spec\bigl(B((q)) \tens{A((q))} (\Z[\zeta_N] \tens{A} (A((q))[X]/(X^N = q)))^{\Gamma \cap \Fix \Lambda_i}\bigr) .
$}
\end{gather*}
If we further suppose that $B$ is flat over $A$, then this last $B((q))$-scheme
\begin{gather*}
\coprod_{\substack{\Lambda_i\ \text{reps of}\\ \HomSurj/\Gamma}}
\Spec\bigl((\Z[\zeta_N] \tens{A} (B((q))[X]/(X^N = q)))^{\Gamma \cap \Fix \Lambda_i}\bigr) .
\end{gather*}
\end{kthm}

\begin{proof}
If $B$ is noetherian and flat over $A$, then $B((q))$ is flat over $A((q))$. Therefore the second assertion results from the first by (\ref{prop:A7.1.3}). To prove the first assertion, we note that it suffices to treat the case $B = A$, for if we denote by $\mathcal{P}_0$ the moduli problem $([\Gamma(N)]/\Gamma)^{A\text{-can}}$ on $\Ell/A$, then we have $\mathcal{P} = \mathcal{P}_0 \tens{A} B$, whence we have (cf.~\ref{sec:4.13}) tautological equalities
\[
\resizebox{\linewidth}{!}{$\displaystyle
\begin{aligned}
\mathcal{P}_{\Tate(q)/B((q))} &= (\mathcal{P}_0)_{\Tate(q)/B((q))} \\
&= ((\mathcal{P}_0)_{\Tate(q)/A((q))}) \tens{A((q))} B((q)) .
\end{aligned}
$}
\]
Finally, the case $B = A$ follows immediately from \eqref{cor:8.11.9} applied to $\Gamma$ acting on $[\Gamma(N)]^{A\text{-can}}$, the fact that $\Tate(q)[N]/A((q))$ is $\mathbf{T}[N] \tens{\Z[q,q^{-1}]} A((q))$, and \eqref{thm:10.5.1} applied in the case $B = A((q))$.
\end{proof}

\begin{kcor}\label{cor:10.8.3}
Let $H$ be a subgroup of $(\Z/N\Z)^\times$, $A$ the subring $\Z[\zeta_N]^H$ of $\Z[\zeta_N]$, and $B$ any flat noetherian $A$-algebra.

Let $\Gamma \subset \GL(2,\Z/N\Z)$ be a subgroup such that $\det(\Gamma) \subset H$, and denote by $\mathcal{P}$ the moduli problem
\[
\mathcal{P} = ([\Gamma(N)]/\Gamma)^{A\text{-can}} \tens{A} B \quad\text{on } \Ell/B \; .
\]
Then we have a canonical isomorphism of $B((q))$-schemes
\begin{gather*}
(\mathcal{P}_{\Tate(q)/B((q))})/\pm 1 \simeq \\
\coprod_{\substack{\Lambda_i\ \text{reps of}\\ \HomSurj/\pm\Gamma}}
\Spec\bigl((\Z[\zeta_N] \tens{A} (B((q))[X]/(X^N = q)))^{(\pm\Gamma) \cap \Fix \Lambda_i}\bigr) .
\end{gather*}
If either $\Gamma$ contains $-1$, or if $-1$ acts without fixed points on the set $\HomSurj/\Gamma$ of cusp-labels for $\Gamma$, then for any $\Lambda \in \HomSurj((\Z/N\Z)^2, \Z/N\Z)$, we have
\end{kcor}

\begin{proof}
Apply the previous theorem to the group $\pm\Gamma$.

If $-1$ has no fixed points acting on $\HomSurj/\Gamma$, then for any element $\gamma \in \Gamma$, we claim $-\gamma$ fixes no $\Lambda \in \HomSurj$. If this were false, then $\Lambda\gamma = -\Lambda$ for some $\gamma \in \Gamma$, $\Lambda \in \HomSurj$, and then $-1$ fixes the coset $\Lambda\Gamma$, contrary to hypothesis.
\end{proof}

\section{Applications to the four basic moduli problems}
\label{sec:10.9}

\begin{kthm}\label{thm:10.9.1}
Let $N \ge 1$ be an integer, $\Gamma \subset \SL(2, \Z/N\Z)$ a subgroup. Denote by $\mathcal{P}$ the moduli problem
\[
([\Gamma(N)]/\Gamma)^{\Z[\zeta_N]\text{-can}} \quad\text{on } (\mathrm{Ell}/\Z[\zeta_N]) ,
\]
and by $\overline{\mathbf{M}}(\mathcal{P})$ its compactified coarse moduli scheme. Then
\begin{enumerate}[label=(\arabic*)]
\item The cusps of $\overline{\mathbf{M}}(\mathcal{P})$ are the disjoint union of $n(\Gamma)$ sections of $\overline{\mathbf{M}}(\mathcal{P})$ over $\Z[\zeta_N]$, where $n(\Gamma)$ is the integer
\[
n(\Gamma) = \#(\HomSurj((\Z/N\Z)^2, \Z/N\Z)/\pm\Gamma) .
\]
\item There exists an open neighborhood $V$ of the cusps
\[
\cusps(\mathcal{P}) \subset V \subset \overline{\mathbf{M}}(\mathcal{P})
\]
which is smooth over $\Z[\zeta_N]$.
\item Pick a set of representatives $\Lambda_i$ for the space $\HomSurj/\pm\Gamma$, and for each representative let $e_i$ be the order of the group $(\pm\Gamma)\cap(\Fix\Lambda_i)$, $f_i = N/e_i$. Then the formal completion of $\overline{\mathbf{M}}(\mathcal{P})$ along its cusps is the $\Z[\zeta_N]$-formal scheme
\[
\coprod_{\Lambda_i \in \HomSurj/\pm\Gamma} \Spf\bigl(\Z[\zeta_N][[q^{1/f_i}]]\bigr) .
\]
\end{enumerate}
\end{kthm}

\begin{proof}
The assertions are all Zariski-local on $\Z[\zeta_N]$, so we may choose an auxiliary prime number $\ell$, and consider instead $\mathcal{P} \tens{\Z[\zeta_N]} \Z[\zeta_N, 1/\ell]$ over the ring $\Z[\zeta_N, 1/\ell]$. By \eqref{thm:8.11.10} and \eqref{thm:10.8.2}, the formal completion of $\overline{\mathbf{M}}(\mathcal{P})[1/\ell]$ along its cusps is the normalization of $\Z[\zeta_N, 1/\ell][[q]]$ in
\[
\coprod_{\Lambda_i \in \HomSurj/\pm\Gamma} \Spec\bigl((\Z[\zeta_N, 1/\ell]((q))[X]/(X^N = q))^{(\pm\Gamma)\cap\Fix(\Lambda_i)}\bigr) .
\]
Because $\pm\Gamma$ lies inside $\SL$, the intersection $(\pm\Gamma)\cap\Fix(\Lambda_i)$ is, via the basis $k_{\Lambda_i}$, $l_{\Lambda_i}$ of $(\Z/N\Z)^2$ (cf.~\ref{prop:10.2.4}) isomorphic to a subgroup of the semi-Borel $\begin{pmatrix} * & * \\ 0 & 1 \end{pmatrix}$, all of whose elements have determinant one. Therefore $(\pm\Gamma)\cap\Fix(\Lambda_i)$ is a subgroup of the upper unipotent group $\begin{pmatrix} 1 & * \\ 0 & 1 \end{pmatrix}$, with respect to the basis $k_{\Lambda_i}$, $l_{\Lambda_i}$. As $\begin{pmatrix} 1 & * \\ 0 & 1 \end{pmatrix}$ is cyclic of order $N$, $(\pm\Gamma)\cap(\Fix\Lambda)$ is the cyclic subgroup generated by the element
\[
\begin{pmatrix} 1 & f_i \\ 0 & 1 \end{pmatrix} ,
\]
operating on the ring $\Z[\zeta_N, 1/\ell]((q))[X]/(X^N = q)$ by
\[
\begin{cases}
\zeta_N \mapsto \zeta_N \\
X \mapsto \zeta_N^{f_i} X .
\end{cases}
\]
Because $(\zeta_N)^{f_i}$ is a primitive $e_i$'th root of unity, the subring of invariants is the subring of polynomials in $X^{e_i} = q^{1/f_i}$. Thus the formal completion of $\overline{\mathbf{M}}(\mathcal{P}) \tens{\Z[\zeta_N]} \Z[\zeta_N, 1/\ell]$ along the cusps is
\[
\coprod_{\Lambda_i \in \HomSurj/\pm\Gamma} \Spec\bigl(\Z[\zeta_N, 1/\ell][[q^{1/f_i}]]\bigr) ,
\]
which is formally smooth over $\Z[\zeta_N]$, and which exhibits the cusps as the asserted number of disjoint sections.
\end{proof}

\begin{kcor}\label{cor:10.9.2}
Suppose that $\Gamma \subset \SL(2,\Z/N\Z)$, and that $K \ge 1$ is an integer such that the moduli problem $\mathcal{P} = ([\Gamma(N)]/\Gamma)^{\mathrm{can}}$ on $(\mathrm{Ell}/\Z[\zeta_N])$ becomes representable when we invert $K$, i.e., $\mathcal{P}[1/K]$ is representable on $(\mathrm{Ell}/\Z[\zeta_N, 1/K])$. Then
\begin{enumerate}[label=(\arabic*)]
\item The compactified moduli scheme $\overline{\mathbf{M}}(\mathcal{P})[1/K]$ is smooth over $\Z[\zeta_N, 1/K]$ except possibly at supersingular points in characteristic $p$ with $p$ a divisor of $N$, $p \nmid K$.
\item The scheme $\overline{\mathbf{M}}(\mathcal{P})[1/NK]$ is a proper smooth curve over $\Z[\zeta_N, 1/NK]$ with geometrically connected fibers.
\end{enumerate}
\end{kcor}

\begin{proof}
We first prove (2). The moduli problem $[\Gamma(N)] \tens{\Z} \Z[1/N]$ is finite etale galois over $(\mathrm{Ell}/\Z[1/N])$, with galois group $\GL(2, \Z/N\Z)$, and finite etale galois over $[\Z[\zeta_N, 1/N]]$ with galois group $\SL(2, \Z/N\Z)$. Therefore $([\Gamma(N)]/\Gamma) \tens{\Z} \Z[1/N]$ is finite etale over $[\Z[\zeta_N, 1/N]]$. By (\ref{cor:9.1.9}), $\mathcal{P}[1/N]$ is finite etale over $(\mathrm{Ell}/\Z[\zeta_N, 1/N])$, so in particular it is smooth over $\Z[\zeta_N, 1/N]$ of relative dimension one. Because $\mathcal{P}[1/K]$ is representable, this means that $\mathbf{M}(\mathcal{P})[1/NK]$ is a smooth curve over $\Z[\zeta_N, 1/NK]$. We already know that $\overline{\mathbf{M}}(\mathcal{P})$ is smooth over $\Z[\zeta_N]$ near its cusps. Therefore $\overline{\mathbf{M}}(\mathcal{P}) \tens{\Z[\zeta_N]} \Z[\zeta_N, 1/NK]$ is a proper smooth curve over $\Z[\zeta_N, 1/NK]$. By the connectedness theorem, it suffices to show that its geometric generic fiber is connected, i.e., to study the situation after the base change $\Z[\zeta_N] \hookrightarrow \mathbb{C}$ defined by, say $\zeta_N \mapsto \exp(2\pi i/N)$. To show that the smooth $\mathbb{C}$-curve $\overline{\mathbf{M}}(\mathcal{P}) \tens{\Z[\zeta_N]} \mathbb{C}$ is connected, it suffices to show that the complement of the finitely many cusps, $\mathbf{M}(\mathcal{P}) \tens{\Z[\zeta_N]} \mathbb{C}$, is connected.
\par
But this is standard, because the underlying complex manifold to $\mathbf{M}(\mathcal{P}) \tens{\Z[\zeta_N]} \mathbb{C}$ is isomorphic to the quotient of the upper half plane by the subgroup $\widetilde{\Gamma} \subset \SL(2,\Z)$ which is the complete inverse image of $\Gamma$ by reduction $\bmod N$.
\par
We now prove (1). Let $p$ be any prime such that $p|N$, $p \nmid K$. Choose an algebraic closure $k$ of $\mathbb{F}_p$, and a homomorphism $\Z[\zeta_N] \to k$.
\par
Let $\mathfrak{p}$ denote the corresponding place of $\Z[\zeta_N]$, and denote by $R$ the completion of the strict henselization of $\Z[\zeta_N]$ at $\mathfrak{p}$, (i.e., $R$ ``is'' the $p$-adic completion of $\Z[\zeta_N]$, with its residue field extended to $k$). Let $E/k$ be an ordinary elliptic curve, $E/R[[T]]$ its universal formal deformation (for deformations to artin local $R$-algebras with residue field $k$). We are claiming that the finite $R[[T]]$-scheme
\[
\mathcal{P}_{E/R[[T]]} ,
\]
viewed as formal scheme over $R$, is formally smooth over $R$. We have proven that $\overline{\mathbf{M}}(\mathcal{P})$ is smooth over $\Z[\zeta_N]$ near the cusps, hence at some ordinary points in every fiber. Therefore there exists an $E/k$ which is ordinary, and such that $\mathcal{P}_{E/R[[T]]}$ is formally smooth over $R$.
\par
We will now see that the isomorphism class of the formal $R$-scheme $\mathcal{P}_{E/R[[T]]}$ is independent of the particular choice of ordinary $E/k$. Because $\mathcal{P}$ is the quotient of $[\Gamma(N)]^{\mathrm{can}}$ by $\Gamma$, we have (\ref{cor:7.1.5})
\[
\mathcal{P}_{E/R[[T]]} = ([\Gamma(N)]^{\mathrm{can}}_{E/R[[T]]})/\Gamma .
\]
By the Serre-Tate theorem, (\ref{thm:2.9.1}), the deformation theory of any elliptic curve $E/k$ to an artin local ring with residue field $k$ is equal to the deformation theory of its $N$-divisible group, given with all the $e_{N^\nu}$-pairings, to that same artin local ring, for any integer $N$ divisible by $p$. Because $E/k$ is ordinary, and $k$ is algebraically closed, there exists an isomorphism
\[
E[N^\infty] \simeq \mu_{N^\infty} \times (\mathbb{Q}/\Z)_{N^\infty\text{-torsion}} ,
\]
compatible with $e_{N^\nu}$-pairings. Therefore the formal $R$-scheme
\[
[\Gamma(N)]^{\mathrm{can}}_{E/R[[T]]}
\]
prorepresents the functor on artin local $R$-algebras with residue field $k$
\[
A \mapsto \left\{
\begin{aligned}
&\text{deformations $G$ to $A$ of $\mu_{N^\infty} \times (\mathbb{Q}/\Z)_{N^\infty\text{-torsion}}$} \\
&\text{with its $e_{N^\nu}$-pairings, plus a $(\Z/N\Z)^2$-generator} \\
&\text{$(P,Q)$ of $G[N]$ with $e_N(P,Q) = \zeta_N$.}
\end{aligned}
\right.
\]
The action of $\Gamma$ is by moving $(P,Q)$. Thus both the formal $R$-scheme $[\Gamma(N)]^{\mathrm{can}}_{E/R[[T]]}$ and the action of $\Gamma$ upon it, are independent of the choice of ordinary $E/k$. Therefore the quotient $\mathcal{P}_{E/R[[T]]}$ is, as formal $R$-scheme, independent of the choice of ordinary $E/k$.
\end{proof}

\kmitem\label{item:10.9.3}Summarizing Table

\begin{center}
{\scshape Application to $[\Gamma(N)]^{\mathrm{can}}$, $[\mathrm{bal.}\ \Gamma_1(N)]^{\mathrm{can}}$}
\end{center}

\begin{center}
\resizebox{\linewidth}{!}{%
\setlength{\tabcolsep}{4pt}%
\renewcommand{\arraystretch}{1.25}%
\begin{tabular}{@{}c@{\quad}l|c|c@{}}
\multicolumn{2}{@{}c|}{\begin{tabular}{@{}c@{}} \textsc{Moduli Pblm} \\ \textsc{over } $\Z[\zeta_N]$ \end{tabular}} &
\begin{tabular}{@{}c@{}} \textsc{Representable} \\ \textsc{over } $\Z[\zeta_N, 1/K]$ \\ \text{for } $K = {?}$ \end{tabular} &
\begin{tabular}{@{}c@{}} $\overline{\mathbf{M}}(\mathcal{P})$ \textsc{Proper Smooth} \\ \textsc{over } $\Z[\zeta_N, 1/K]$ \\ \textsc{outside s/s points} \\ \textsc{in char } $p = {?}$ \end{tabular} \\
\hline\hline
\multirow{4}{*}{$[\Gamma(N)]^{\mathrm{can}}$} &
\multicolumn{1}{@{}l}{$\left\{\begin{array}{@{}l@{}}
N = 1 \text{ or } 2 \\
N = p^k \ge 3 \\
N = 2p^k, p \ne 2 \\
\parbox{0.32\linewidth}{$N$ divisible by two rel.\ prime integers both $\ge 3$}
\end{array}\right.$} &
\begin{tabular}{c}
\textsc{Never} \\ \hline
$p$ \\ \hline
$p$ \\ \hline
$1$
\end{tabular} &
\begin{tabular}{c}
\textsc{(Not Applicable)} \\ \hline
\textsc{None Left} \\ \hline
$2$ \\ \hline
all $p|N$
\end{tabular} \\
\hline\hline
\multirow{6}{*}{$[\mathrm{bal.}\ \Gamma_1(N)]^{\mathrm{can}}$} &
\multicolumn{1}{@{}l}{$\left\{\begin{array}{@{}l@{}}
N = 1,2,3 \\
N = p^k \ge 4 \\
N = 2p^k, p \text{ odd}, p^k \ge 4 \\
N = 3p^k, p \ne 3, p^k \ge 4 \\
N = 6 \\
\parbox{0.32\linewidth}{$N$ divisible by two rel.\ prime integers both $\ge 4$}
\end{array}\right.$} &
\begin{tabular}{c}
\textsc{Never} \\ \hline
$p$ \\ \hline
$p$ \\ \hline
$p$ \\ \hline
$6$ \\ \hline
$1$
\end{tabular} &
\begin{tabular}{c}
\textsc{(Not Applicable)} \\ \hline
\textsc{None Left} \\ \hline
$2$ \\ \hline
$3$ \\ \hline
\textsc{None Left} \\ \hline
all $p|N$
\end{tabular} \\
\hline\hline
\end{tabular}%
}
\end{center}

\begin{kthm}\label{thm:10.9.4}
Let $N_1$ and $N_2$ be integers $\ge 1$ which are relatively prime: $(N_1, N_2) = 1$. Put $N = N_1N_2$, and let $\Gamma \subset \GL(2, \Z/N\Z)$ be a subgroup which satisfies
\[
\det(\Gamma) \equiv 1 \bmod N_1 \; .
\]
Let $A$ denote the subring $\Z[\zeta_N]^{\det(\Gamma)}$ of $\Z[\zeta_N]$, $\mathcal{P}$ the moduli problem
\[
([\Gamma(N)]/\Gamma)^{A\text{-can}} \quad\text{on } (\mathrm{Ell}/A) \; .
\]
Denote by $\Gamma_0$ the intersection $\Gamma \cap \SL(2, \Z/N\Z)$, and by $\mathcal{P}_0$ the moduli problem
\[
([\Gamma(N)]/\Gamma_0)^{\Z[\zeta_N]\text{-can}} \quad\text{on } (\mathrm{Ell}/\Z[\zeta_N]) \; .
\]
Then
\begin{enumerate}[label=(\arabic*)]
\item We have an isomorphism of moduli problems over $\Z[\zeta_N, \allowbreak 1/N_2]$
\[
\mathcal{P}[1/N_2] \tens{A[1/N_2]} \Z[\zeta_N, 1/N_2] \xrightarrow{\sim} \mathcal{P}_0[1/N_2] ,
\]
\item This isomorphism induces an isomorphism of compactified moduli schemes
\[
\overline{\mathbf{M}}(\mathcal{P}) \tens{A} \Z[\zeta_N, 1/N_2] \xrightarrow{\sim} \overline{\mathbf{M}}(\mathcal{P}_0)[1/N_2] .
\]
\item The scheme $\overline{\mathbf{M}}(\mathcal{P})[1/N_2]$ is smooth over $A[1/N_2]$ in a neighborhood of the cusps, and its cusps are a finite etale $A[1/N_2]$ scheme of degree $n(\Gamma_0)$,
\[
n(\Gamma_0) = \#(\HomSurj((\Z/N\Z)^2, \Z/N\Z)/\pm\Gamma_0) .
\]
\end{enumerate}
\end{kthm}

\begin{proof}
Let $H = \det(\Gamma)$ inside $(\Z/N\Z)^\times$, and let $H_2$ be the image of $H$ in $(\Z/N_2\Z)^\times$. Because $\det H \equiv 1(N_1)$, we have $H \xrightarrow{\sim} H_2$ operating on
\[
\Z[\zeta_N] = \Z[\zeta_{N_1}] \tens{\Z} \Z[\zeta_{N_2}]
\]
through the second factor. Therefore, we have
\[
A = \Z[\zeta_N]^{\det(\Gamma)} = \Z[\zeta_{N_1}] \tens{\Z} (\Z[\zeta_{N_2}])^{H_2} .
\]
But $\Z[\zeta_{N_2}, 1/N_2]$ is finite etale over $\Z[1/N_2]$, as is its subring of $H_2$-invariants. Therefore $\Z[\zeta_N, 1/N_2]$ is finite etale over $A[1/N_2]$.
\par
We now apply \eqref{prop:9.2.2} to the situation $G = \Gamma/\Gamma_0$ acting on $([\Gamma(N)]/\Gamma_0)[1/N]$, covering its free action on $\Z[\zeta_N, 1/N_2]$. This proves (1), and (2) follows from (1) and the fact that $\Z[\zeta_N, 1/N_2]$ is etale over $A[1/N_2]$ (cf.~\ref{prop:8.6.7}, (3)). The assertions of (3) are f.p.p.f.\ local on the base, and we have already proven, thanks to (2) and \eqref{cor:10.9.2}, that they become true after the finite etale base change $A[1/N_2] \to \Z[\zeta_N, 1/N_2]$.
\end{proof}

\begin{kcor}\label{cor:10.9.5}
Hypotheses and notations as in \eqref{thm:10.9.4}, suppose further that for some integer $K \ge 1$ the moduli problem $([\Gamma(N)]/\Gamma)[1/K]$ on $(\mathrm{Ell}/\Z[1/K])$ is representable. Let $\mathcal{P}$ denote, as before, the moduli problem $([\Gamma(N)]/\Gamma)^{A\text{-can}}$ on $(\mathrm{Ell}/A)$. Then $\mathcal{P}[1/N_2K]$ is representable, by \eqref{prop:9.1.7}, and we have the following results.
\begin{enumerate}[label=(\arabic*)]
\item The compactified moduli scheme $\overline{\mathbf{M}}(\mathcal{P})[1/N_2K]$ is smooth over $A[1/N_2K]$ except possibly at supersingular points in characteristic $p$ with $p$ a divisor of $N_1$, $p \nmid KN_2$.
\item The scheme $\overline{\mathbf{M}}(\mathcal{P})[1/NK]$ is a proper smooth curve over $A[1/NK]$ with geometrically connected fibers, in which the scheme of cusps is finite etale over $A[1/NK]$.
\end{enumerate}
\end{kcor}

\begin{proof}
The statements are f.p.p.f.\ local on the base, and we know already they become true after the finite etale base changes $A[1/N_2K] \to \Z[\zeta_N, 1/N_2K]$ and $A[1/NK] \to \Z[\zeta_N, 1/NK]$ respectively, because the isomorphism (1) of \eqref{thm:10.9.4} shows that $\mathcal{P}_0[1/N_2K]$ is representable, so that \eqref{cor:10.9.2} applies to it.
\end{proof}

\kmitem\label{item:10.9.6}Summarizing Table

\begin{center}
{\scshape Short Table of Applications}
\end{center}

\begin{center}
\resizebox{\linewidth}{!}{%
\setlength{\tabcolsep}{4pt}%
\renewcommand{\arraystretch}{1.3}%
\begin{tabular}{@{}c|c|c|c@{}}
\hline
\begin{tabular}{@{}c@{}} \textsc{Moduli Pblm} $\mathcal{P}$ \\ \textsc{of form} $([\Gamma(N_1N_2)]/\Gamma)^{A\text{-can}}$ \\ \text{with } $(N_1, N_2) = 1$, \\ $\det\Gamma \equiv 1 \bmod N_1$ \end{tabular} &
$A = {?}$ &
\begin{tabular}{@{}c@{}} \textsc{Rep Over} \\ $A[1/K] = {?}$ \end{tabular} &
\begin{tabular}{@{}c@{}} $\overline{\mathbf{M}}(\mathcal{P})[1/N_2K]$ \textsc{Proper} \\ \textsc{Smooth Over} \\ $A[1/N_2K]$ \textsc{outside} \\ \textsc{s/s points in} \\ \textsc{char} $p|N_1$, \\ $A[1/N_2K] = {?}$ \end{tabular} \\
\hline\hline
\begin{tabular}{@{}c@{}} $([\Gamma(N_1)]^{\Z[\zeta_{N_1}]\text{-can}}, [\Gamma_0(N_2)])$ \\ \text{with } $N_1 \ge 3$ \end{tabular} &
$\Z[\zeta_{N_1}]$ &
$\Z[\zeta_{N_1}, 1/N_1]$ &
$\Z[\zeta_{N_1}, 1/N_1N_2]$ \\
\hline
\begin{tabular}{@{}c@{}} $([\mathrm{bal.}\ \Gamma_1(N_1)]^{\Z[\zeta_{N_1}]\text{-can}}, [\Gamma_0(N_2)])$ \\ \text{with } $N_1 \ge 4$ \end{tabular} &
$\Z[\zeta_{N_1}]$ &
$\Z[\zeta_{N_1}, 1/N_1]$ &
$\Z[\zeta_{N_1}, 1/N_1N_2]$ \\
\hline
\begin{tabular}{@{}c@{}} $([\Gamma_1(N_1)], [\Gamma_0(N_2)])$ \\ \text{with } $N_1 \ge 4$ \end{tabular} &
$\Z$ &
$\Z[1/N_1]$ &
$\Z[1/N_1N_2]$ \\
\hline
$[\Gamma_1(N_1)]$ \text{with } $N_1 \ge 4$ &
$\Z$ &
$\Z[1/N_1]$ &
$\Z[1/N_1]$ \\
\hline
\begin{tabular}{@{}c@{}} $([\mathrm{bal.}\ \Gamma_1(N_1)]^{\Z[\zeta_{N_1}]\text{-can}}, [\Gamma_1(N_2)])$ \\ \text{with } $N_2 \ge 4$ \end{tabular} &
$\Z[\zeta_{N_1}]$ &
$\Z[\zeta_{N_1}][1/N_2]$ &
$\Z[\zeta_{N_1}][1/N_2]$ \\
\hline
\end{tabular}%
}
\end{center}

\begin{kthm}\label{thm:10.9.7}
Let $\Gamma \subset \GL(2, \Z/N\Z)$ be a subgroup, $A = (\Z[\zeta_N])^{\det(\Gamma)}$ the corresponding subring of $\Z[\zeta_N]$, and denote by $\discrim(\Gamma)$ the product of those prime numbers $p$ such that the extension $\Z[\zeta_N]/A$ is ramified at some $p$-adic place of $A$ (so $\discrim(\Gamma)$ is always a divisor of $N$).
\par
Denote by $\mathcal{P}$ the moduli problem
\[
([\Gamma(N)]/\Gamma)^{A\text{-can}} \quad\text{on } (\mathrm{Ell}/A) \; .
\]
The compactified coarse moduli scheme $\overline{\mathbf{M}}(\mathcal{P})[1/\discrim(\Gamma)]$ is smooth over $A[1/\discrim(\Gamma)]$ except possibly [at points of $\mathbf{M}(\mathcal{P})$ at which either $j(j-1728)$ is not invertible, or]{\renewcommand{\thefootnote}{*}\footnote{The phrase in brackets can be deleted, see Notes Added in Proof.}} at points which are supersingular points in characteristic $p$ dividing $N$.
\end{kthm}

\begin{proof}
By the descent mechanism already employed \eqref{thm:10.9.4} we may replace $\Gamma$ by $\Gamma \cap \SL(2, \Z/N\Z)$, and $A$ by $\Z[\zeta_N]$ itself. Let $k$ be the algebraic closure of a finite field, $j_0 \in k$ a number with $j_0 \ne 0, 1728$, $\Z[\zeta_N] \to k$ a homomorphism, $\mathfrak{p}$ the place it defines, and $R$ the completion of the strict henselization of $\Z[\zeta_N]$ at $\mathfrak{p}$.
\par
Let $E/k$ be an elliptic curve with $j(E/k) = j_0$ different from $0, 1728$. Let $E/R[[T]]$ be its universal formal deformation. The formal fiber of $\mathbf{M}(\mathcal{P}) \to \Spec(\Z[\zeta_N][j])$ over $E/R[[T]]$ is \eqref{prop:8.2.3} the formal $R$-scheme
\[
\mathcal{P}_{E/R[[T]]}/\Aut(E/k) .
\]
We must show that it is formally smooth over $R$ except if $E/k$ is supersingular in characteristic $p|N$. Because $\Aut(E/k) = \pm 1$, we may use (\ref{cor:7.1.5}) to rewrite this
\[
([\Gamma(N)]^{\mathrm{can}}_{E/R[[T]]})/\pm \Gamma_0 .
\]
Because we already know smoothness along some neighborhood of the cusps, it suffices to show that for each fixed situation $\Z[\zeta_N] \to k$, this formal $R$-scheme is independent of the particular choice of $E/k$, provided that $E/k$ is ordinary in case $\mathrm{char}(k)|N$. When $\mathrm{char}(k)|N$, this follows from Serre-Tate theory, exactly as in the proof of \eqref{cor:10.9.2}. If $N$ is invertible in $k$, then $E[N]/R[[T]]$ is isomorphic to the extension of scalars of the group over $\Z$
\[
\mu_N \times \Z/N\Z
\]
with its canonical $e_N$!
\end{proof}

\section{Detailed analysis at a prime $p$; balanced subgroups}
\label{sec:10.10}

\kmitem\label{item:10.10.1}\textbf{Definitions.} Let $p$ be a prime number, $n \ge 1$ an integer, and $\Gamma \subset \GL(2, \Z/p^n\Z)$ a subgroup. We say that $\Gamma$ satisfies $(*)$ if
\begin{equation}\label{eq:10.10.1.1}
\begin{aligned}
(*) \quad & \text{either } p \ne 2, \text{ or } p = 2 \text{ and } n = 1, \text{ or } p = 2, n \ge 2, \\
& \text{and } \det(\Gamma) \equiv 1 \bmod 4 .
\end{aligned}
\end{equation}
\par
We say that $\Gamma$ is \emph{balanced} if for every element
\[
\Lambda \in \HomSurj((\Z/p^n\Z)^2, \Z/p^n\Z) ,
\]
the subgroup $\Gamma \cap \Fix(\Lambda)$ of $\Gamma$ has exactly the same group of determinants as $\Gamma$ itself, i.e., if
\begin{equation}\label{eq:10.10.1.2}
\det(\Gamma \cap \Fix(\Lambda)) = \det(\Gamma) \quad\text{for all } \Lambda .
\end{equation}

\par\addvspace{6pt plus 2pt minus 1pt}\noindent\phantomsection\label{item:10.10.1.3}(10.10.1.3)\enspace We say that $\Gamma$ is ``\emph{balanced at $\Lambda$}'' if the above equality of determinant groups holds for the particular element $\Lambda \in \HomSurj$.
\par
Clearly for given $\Gamma$, the question of whether $\Gamma$ is balanced at $\Lambda$ depends only on the image of $\Lambda$ in the space
\[
(\Z/p^n\Z)^\times \backslash \HomSurj/\Gamma
\]
of component-labels for $\Gamma$. We say that $\Gamma$ is \emph{partly balanced} if there exists some component-label at which it is balanced, and we say the corresponding component label is \emph{$\Gamma$-balanced}.

\kmitem\label{item:10.10.2}
Now let $N \ge 1$ be an integer prime to $p$, and consider a subgroup
\[
\Gamma \subset \GL(2, \Z/Np^n\Z)
\]
which is a \emph{product} $\Gamma = \Gamma_1 \times \Gamma_2$ of subgroups
\begin{align*}
\Gamma_1 &\subset \GL(2, \Z/p^n\Z) \\
\Gamma_2 &\subset \GL(2, \Z/N\Z)
\end{align*}
via the canonical isomorphism
\[
\GL(2, \Z/Np^n\Z) \xrightarrow{\sim} \GL(2, \Z/p^n\Z) \times \GL(2, \Z/N\Z)
\]
of the Chinese Remainder Theorem. Recall (\ref{lem:10.7.3}) that this isomorphism induces isomorphisms
\[
\resizebox{0.95\linewidth}{!}{$\displaystyle
\left\{
\begin{aligned}
\HomSurj &\xrightarrow{\sim} \HomSurj \times \HomSurj \\
\Lambda &\mapsto (\Lambda_1, \Lambda_2) \\
\text{cusp labels for } \Gamma &\xrightarrow{\sim} (\text{cusp labels for } \Gamma_1) \times (\text{cusp labels for } \Gamma_2) \\
\text{compt.\ labels for } \Gamma &\xrightarrow{\sim} (\text{compt.\ labels for } \Gamma_1) \times (\text{compt.\ labels for } \Gamma_2) .
\end{aligned}
\right.
$}
\]
We say that a component label for $\Gamma$ is \emph{$\Gamma_1$-balanced} if its first component via the above isomorphism is a $\Gamma_1$-balanced component label for $\Gamma_1$. We say that $\Gamma$ is \emph{balanced at $p$} if every component label for $\Gamma$ is $\Gamma_1$-balanced, i.e., if $\Gamma_1$ is balanced.

\begin{kthm}\label{thm:10.10.3}
In the above product situation $\Gamma = \Gamma_1 \times \Gamma_2$, $\Gamma_1 \subset \GL(2, \Z/p^n\Z)$, $\Gamma_2 \subset \GL(2, \Z/N\Z)$, $N$ prime to $p$, suppose that
\begin{enumerate}[label=(\alph*)]
\item $\Gamma_1$ satisfies $(*)$ \eqref{eq:10.10.1.1}
\item either both $\Gamma_1$ and $\Gamma_2$ contain $-1$, or $-1$ operates without fixed points on the space of cusp labels for $\Gamma$.
\end{enumerate}
Let
\begin{align*}
A &= (\Z[\zeta_{p^n}])^{\det(\Gamma_1)} \subset \Z[\zeta_{p^n}] , \\
B &= (\Z[\zeta_N])^{\det(\Gamma_2)} \subset \Z[\zeta_N] ,
\end{align*}
so that
\[
A \tens{\Z} B = (\Z[\zeta_{Np^n}])^{\det(\Gamma)} \subset \Z[\zeta_{Np^n}] .
\]
Let $\mathcal{P}$ denote the moduli problem
\[
([\Gamma(Np^n)]/\Gamma)^{A \tens{\Z} B\text{-can}} \quad\text{on } (\mathrm{Ell}/A \tens{\Z} B) .
\]
Then
\begin{enumerate}[label=(\arabic*)]
\item The compactified coarse moduli scheme $\overline{\mathbf{M}}(\mathcal{P})[1/N]$ is \emph{regular} in some neighborhood of the cusps. In fact it is regular except possibly at points at which either $j(j-1728)$ fails to be invertible, or at supersingular points in characteristic $p$ ($p = \text{the given } p$).
\item If $\mathcal{P}[1/N]$ is representable, then $\overline{\mathbf{M}}(\mathcal{P})[1/N]$ is regular outside of the supersingular points in characteristic $p$.
\item The completion $\widehat{\mathrm{Cusps}}(\mathcal{P}[1/N])$ of $\overline{\mathbf{M}}(\mathcal{P})[1/N]$ along its cusps is a disjoint union of formal schemes indexed by
\[
\pm 1 \backslash \HomSurj((\Z/Np^n\Z)^2, \Z/Np^n\Z)/\Gamma .
\]
For each $\Lambda$ in a set of representatives in $\HomSurj$ of this indexing set, let $\Lambda_1$ and $\Lambda_2$ be the components of $\Lambda$, and denote by $A_\Lambda$ and $B_\Lambda$ the rings
\begin{align*}
A \subset A_\Lambda &= (\Z[\zeta_{p^n}])^{\det(\Gamma_1 \cap \Fix(\Lambda_1))} \\
B \subset B_\Lambda &= (\Z[\zeta_N])^{\det(\Gamma_2 \cap \Fix(\Lambda_2))}
\end{align*}
so that $A_\Lambda \tens{\Z} B_\Lambda$ is the subring of invariants of $\det(\Gamma \cap \Fix(\Lambda))$ in $\Z[\zeta_{Np^n}]$. The $\Lambda$-piece of $\widehat{\mathrm{Cusps}}(\mathcal{P}[1/N])$ is naturally a scheme over $A_\Lambda \tens{\Z} B_\Lambda[1/N]$, isomorphic to the formal completion of a smooth $A_\Lambda \tens{\Z} B_\Lambda[1/N]$-scheme of relative dimension one along a closed subscheme of $\mathrm{Cusps}(\mathcal{P})[1/N]$ which is finite etale over $A_\Lambda \tens{\Z} B_\Lambda[1/N]$.
\item If $\Lambda$ is $\Gamma_1$-balanced, then $A = A_\Lambda$, and $\overline{\mathbf{M}}(\mathcal{P})[1/N]$ is smooth over $A \tens{\Z} B[1/N]$ in some neighborhood of the $\Lambda$-cusps. In particular, if $\Gamma_1$ is partly balanced, the ``\emph{ouvert de lissit\'e}'' of $\overline{\mathbf{M}}(\mathcal{P})[1/N]$ over $A \tens{\Z} B[1/N]$ meets every fiber.
\item If $\Gamma_1$ is balanced, then $\overline{\mathbf{M}}(\mathcal{P})[1/N]$ is smooth over $A \tens{\Z} B[1/N]$ except possibly [at points at which either $j(j-1728)$ is not invertible, or]{\renewcommand{\thefootnote}{*}\footnote{The phrase in brackets can be deleted, see Notes Added in Proof.}} at supersingular points in characteristic $p$. Its scheme of cusps is finite etale over $A \tens{\Z} B[1/N]$.
\item If $\Gamma_1$ is balanced and if $\mathcal{P}[1/N]$ is representable, then $\overline{\mathbf{M}}(\mathcal{P})[1/N]$ is smooth over $A \tens{\Z} B[1/N]$ outside of the supersingular points in characteristic $p$. Its scheme of cusps is finite etale over $A \tens{\Z} B[1/N]$.
\end{enumerate}
\end{kthm}

\begin{proof}
The basic assertion to prove is (3). Admitting it for a moment, let us explain how the others follow. Each $A_\Lambda \tens{\Z} B_\Lambda$ is the full ring of integers in a subfield of $\mathbb{Q}(\zeta_{Np^n})$, so is regular. Therefore the $\Lambda$-piece of $\widehat{\mathrm{Cusps}}(\mathcal{P}[1/N])$ is regular, and so $\overline{\mathbf{M}}(\mathcal{P})[1/N]$ is regular at its cusps. As it is a scheme of finite type over $\Z$, its set of regular points is open. The usual homogeneity argument then gives (1) and (2). Assertions (4), (5) and (6) follow similarly, for $B_\Lambda[1/N]$ is finite etale over $B[1/N]$, so $A \tens{\Z} B_\Lambda[1/N]$ is finite etale over $A \tens{\Z} B[1/N]$. Thus whenever $A = A_\Lambda$, we can replace ``regular'' by ``smooth over $A \tens{\Z} B[1/N]$'' in the above argument.
\par
We now prove (3). By \eqref{cor:10.8.3}, we know that $\widehat{\mathrm{Cusps}}(\mathcal{P}[1/N])$ is the normalization of $\Z[1/N][[q]]$ in the finite $\Z[1/N]((q))$-scheme
\[
\resizebox{0.97\linewidth}{!}{$\displaystyle
\coprod_{\substack{\Lambda\ \text{reps of}\\ \HomSurj/\pm\Gamma}} \Spec\bigl((\Z[\zeta_{p^n}, \zeta_N, 1/N]((q))[X]/(X^{Np^n} = q))^{(\pm\Gamma)\cap\Fix(\Lambda)}\bigr) .
$}
\]
By hypothesis (b), we have (cf.~\ref{cor:10.8.3})
\[
(\pm\Gamma)\cap\Fix\Lambda = \Gamma \cap \Fix\Lambda .
\]
Because normalization commutes with passage to invariants, we have that $\widehat{\mathrm{Cusps}}(\mathcal{P}[1/N])$ is the disjoint union, over the representatives $\Lambda$, of the spectra of the rings
\[
(\Z[\zeta_{p^n}, \zeta_N, 1/N][[q]][X]/(X^{Np^n} = q))^{\Gamma \cap \Fix\Lambda} .
\]
Because $\Z[1/N][[q]]$ is flat over $\Z[1/N][q]$, the above ring may be written as
\[
\Z[1/N][[q]] \tens{\Z[1/N, q]} (\Z[\zeta_{p^n}, \zeta_N, 1/N, q, X]/(X^{Np^n} = q))^{\Gamma \cap \Fix\Lambda} .
\]
Thus we are reduced to showing that the spectrum of
\[
(\Z[\zeta_{p^n}, \zeta_N, 1/N, q, X]/(X^{Np^n} = q))^{\Gamma \cap \Fix\Lambda}
\]
is smooth and one-dimensional over $(\Z[\zeta_{p^n}, \zeta_N, 1/N])^{\det(\Gamma \cap \Fix\Lambda)}$, and that the closed reduced subscheme of it defined set-theoretically by the equation $q = 0$ is finite etale over the same ring.
\par
We next reduce to the case when $\Gamma_2$ lies in $\SL$.
\par
Let $\Gamma_2' = \Gamma_2 \cap \SL(2, \Z/N\Z)$, $\Gamma' = \Gamma_1 \times \Gamma_2'$ inside $\Gamma = \Gamma_1 \times \Gamma_2$. Then
\begin{align*}
\Gamma' \cap \Fix(\Lambda) &= (\Gamma_1 \cap \Fix\Lambda_1) \times (\Gamma_2' \cap \Fix\Lambda_2) \\
&= (\Gamma_1 \cap \Fix\Lambda_1) \times (\Gamma_2 \cap \Fix\Lambda_2 \cap \SL(2, \Z/N\Z)) ,
\end{align*}
so that the quotient of $\Gamma \cap \Fix\Lambda$ by $\Gamma' \cap \Fix\Lambda$ is isomorphic, via the determinant, to the galois group of the finite etale galois covering
\begin{center}
\begin{tikzcd}[row sep=2.5em]
A_\Lambda \tens{\Z} \Z[\zeta_N, 1/N] = (\Z[\zeta_{p^n}, \zeta_N, 1/N])^{\det(\Gamma' \cap \Fix\Lambda)} \arrow[d, dash] \\
A_\Lambda \tens{\Z} B_\Lambda[1/N] = (\Z[\zeta_{p^n}, \zeta_N, 1/N])^{\det(\Gamma \cap \Fix\Lambda)} .
\end{tikzcd}
\end{center}
Therefore the situation we are considering is a descent through this finite etale galois covering of the situation
\[
(\Z[\zeta_{p^n}, \zeta_N, 1/N, q, X]/(X^{Np^n} = q))^{\Gamma' \cap \Fix\Lambda}
\]
with respect to the ring $A_\Lambda[\zeta_N, 1/N]$. This completes the reduction to the case $\Gamma_2 \subset \SL$.
\par
Supposing now that $\Gamma_2 \subset \SL$, consider the product decomposition
\[
\Gamma \cap \Fix\Lambda = (\Gamma_1 \cap \Fix\Lambda_1) \times (\Gamma_2 \cap \Fix\Lambda_2) .
\]
In the standard bases $(k_{\Lambda_1}, l_{\Lambda_1})$ and $(k_{\Lambda_2}, l_{\Lambda_2})$ adapted to $\Lambda_1$ and $\Lambda_2$ respectively (cf.~\ref{prop:10.2.4}) we have matrix realizations
\begin{align*}
\Gamma_1 \cap \Fix\Lambda_1 &\subset \begin{pmatrix} * & * \\ 0 & 1 \end{pmatrix} \subset \GL(2, \Z/p^n\Z) \\
\Gamma_2 \cap \Fix\Lambda_2 &\subset \begin{pmatrix} 1 & * \\ 0 & 1 \end{pmatrix} \subset \GL(2, \Z/N\Z) .
\end{align*}
Let us denote by $e$ the order of the group $\Gamma_2 \cap \Fix\Lambda_2$, $f$ the ratio $N/e$, its index in $\begin{pmatrix} 1 & * \\ 0 & 1 \end{pmatrix}$. Then visibly the $\Gamma_2 \cap \Fix\Lambda_2$ invariants in
\[
\Z[\zeta_{p^n}, \zeta_N, 1/N, q, X]/(X^{Np^n} = q)
\]
are the polynomials in $X^e$.
\par
Introducing the new variable $T = X^e$, and replacing $\zeta_{p^n}$ by its unique $f$'th root which is also a $p^n$-th root of unity, we are reduced to studying the ring of invariants
\[
(\Z[\zeta_{p^n}, \zeta_N, 1/N, q, T]/(T^{fp^n} = q))^{\Gamma_1 \cap \Fix\Lambda_1} ,
\]
where an element
\[
\begin{pmatrix} A & B \\ 0 & 1 \end{pmatrix} \in \Gamma_1 \cap \Fix\Lambda_1 \subset \GL(2, \Z/p^n\Z)
\]
acts by
\[
\left\{
\begin{aligned}
\zeta_{p^n} &\mapsto (\zeta_{p^n})^A \\
\zeta_N &\mapsto \zeta_N \\
q &\mapsto q \\
T &\mapsto T(\zeta_{p^n})^B .
\end{aligned}
\right.
\]
We must show that the ring of invariants is smooth over $A_\Lambda[\zeta_N, 1/N]$, and that the reduced closed subscheme $(q=0)^{\mathrm{red}}$ is finite etale over $A_\Lambda[\zeta_N, 1/N]$.
\par
The situation in question comes by the etale extension of scalars $\Z \to \Z[\zeta_N, 1/N]$ from the analogous situation
\[
(\Z[\zeta_{p^n}, T], \text{above action of } \Gamma_1 \cap \Fix(\Lambda)) .
\]
It is at this point we make use of the hypothesis that $\Gamma_1$ satisfies $(*)$, i.e., that if $p = 2$ and $p^n$ is $4$ or more, then $\det(\Gamma_1) \equiv 1 \bmod 4$. Obviously any subgroup of $\Gamma_1$ also satisfies $(*)$, in particular the group $\Gamma_1 \cap \Fix\Lambda$. We thus conclude the proof by appealing to the following theorem.
\end{proof}

\begin{kthm}\label{thm:10.10.4}
Let $G \subset \begin{pmatrix} * & * \\ 0 & 1 \end{pmatrix} \subset \GL(2, \Z/p^n\Z)$ be a subgroup of the standard upper semi-Borel, and suppose that $G$ satisfies $(*)$ (cf.~\ref{eq:10.10.1.1}). Let us denote
\[
p^a = \#\left(G \cap \begin{pmatrix} 1 & * \\ 0 & 1 \end{pmatrix}\right) .
\]
Then there exists a root of unity
\[
\zeta \in \Z[\zeta_{p^n}]
\]
such that under the standard action of $G$ on $\Z[\zeta_{p^n}, T]$
\[
\begin{pmatrix} A & B \\ 0 & 1 \end{pmatrix} :
\left\{
\begin{aligned}
\zeta_{p^n} &\mapsto (\zeta_{p^n})^A \\
T &\mapsto T(\zeta_{p^n})^B ,
\end{aligned}
\right.
\]
the ring of invariants is the polynomial ring in one variable
\[
(\Z[\zeta_{p^n}, T])^G = (\Z[\zeta_{p^n}])^{\det(G)}[S], \quad S = \zeta T^{p^a} .
\]
\end{kthm}

\begin{proof}
This is simply group-theory. Consider the exact sequence of groups
\[
1 \longrightarrow G \cap \begin{pmatrix} 1 & * \\ 0 & 1 \end{pmatrix} \longrightarrow G \xrightarrow{\det} (\Z/p^n\Z)^\times .
\]
The subgroup $G \cap \begin{pmatrix} 1 & * \\ 0 & 1 \end{pmatrix}$, having order $p^a$, is the cyclic group generated by
\[
\begin{pmatrix} 1 & p^{n-a} \\ 0 & 1 \end{pmatrix} .
\]
If $G = G \cap \begin{pmatrix} 1 & * \\ 0 & 1 \end{pmatrix}$, the assertion is obvious (with $\zeta = 1$). If $\det(G)$ is non-trivial, let $A \in (\Z/p^n\Z)^\times$ be a generator of the cyclic (by $(*)$) group $\det(G)$, and let
\[
\begin{pmatrix} A & B \\ 0 & 1 \end{pmatrix}
\]
be some chosen element of $G$ with determinant $A$. These two elements obviously generate $G$.
\par
The subring of invariants under $G \cap \begin{pmatrix} 1 & * \\ 0 & 1 \end{pmatrix}$ is visibly the ring
\[
\Z[\zeta_{p^n}, T^{p^a}] ,
\]
on which $\begin{pmatrix} A & B \\ 0 & 1 \end{pmatrix}$ operates by
\begin{align*}
\zeta_{p^n} &\mapsto (\zeta_{p^n})^A \\
T^{p^a} &\mapsto T^{p^a}(\zeta_{p^n})^{Bp^a} .
\end{align*}
We will produce a root of unity $\zeta \in \Z[\zeta_{p^n}]$, such that the element
\[
S = \zeta T^{p^a}
\]
is invariant by $\begin{pmatrix} A & B \\ 0 & 1 \end{pmatrix}$, and hence by all of $G$. Once we have done this, then
\begin{align*}
(\Z[\zeta_{p^n}, T])^G &= (\Z[\zeta_{p^n}, T^{p^a}])^{\begin{pmatrix} A & B \\ 0 & 1 \end{pmatrix}} \\
&= (\Z[\zeta_{p^n}, S])^{\begin{pmatrix} A & B \\ 0 & 1 \end{pmatrix}} \\
&= (\Z[\zeta_{p^n}])^{\det(G)}[S] .
\end{align*}
We now search for $\zeta$ of the form $(\zeta_{p^n})^x$, such that $\zeta T^{p^a}$ is invariant by $\begin{pmatrix} A & B \\ 0 & 1 \end{pmatrix}$:
\[
T^{p^a}(\zeta_{p^n})^x \xmapsto{\begin{pmatrix} A & B \\ 0 & 1 \end{pmatrix}} T^{p^a}(\zeta_{p^n})^{Ax}(\zeta_{p^n})^{Bp^a} .
\]
Thus the requirement is that
\[
x \equiv Ax + p^a B \bmod p^n ,
\]
or equivalently
\[
(1 - A)x \equiv p^a B \bmod p^n .
\]
If $A \not\equiv 1 \bmod p$, then $1 - A$ is invertible $\bmod\ p^n$, so we can solve uniquely for $x$.
\par
If $A \equiv 1 \bmod p$, then $\det(G)$ is the subgroup
\[
1 + (p^b) \subset (\Z/p^n\Z)^\times
\]
for some $1 \le b \le n-1$. ($\le n-1$ because we have already dealt with the case $\det(G) = 1$.) If $p = 2$, then $b \ge 2$, because of condition $(*)$. We may choose for $A$ the generator
\[
A = 1 + p^b .
\]
We must show that
\[
p^b x \equiv p^a B \bmod p^n
\]
has a solution. If $B \equiv 0 \bmod p^n$, then $x = 0$ is a solution. If $B \not\equiv 0 \bmod p^n$, we may speak of the integer $\ord_p(B) \le n-1$, and we must show that
\[
\ord_p(B) + a - b \ge 0 ,
\]
given that the matrix
\[
\begin{pmatrix} A & B \\ 0 & 1 \end{pmatrix} = \begin{pmatrix} 1 + p^b & B \\ 0 & 1 \end{pmatrix}
\]
has its $p^{n-b}$'th power in $G \cap \begin{pmatrix} 1 & * \\ 0 & 1 \end{pmatrix}$, which is the cyclic group generated by
\[
\begin{pmatrix} 1 & p^{n-a} \\ 0 & 1 \end{pmatrix} .
\]
For any integer $m \ge 1$, one easily verifies the identity
\[
\begin{pmatrix} A & B \\ 0 & 1 \end{pmatrix}^m = \begin{pmatrix} A^m & B(1 + A + \dots + A^{m-1}) \\ 0 & 1 \end{pmatrix} .
\]
In our situation, this yields the congruence
\[
B(1 + A + \dots + A^{p^{n-b}-1}) \equiv 0 \bmod p^{n-a} .
\]
In terms of $A$, $B \in \Z_p$ lifting $A$, $B \in \Z/p^n\Z$, this gives
\[
B\left(\frac{A^{p^{n-b}}-1}{A-1}\right) \equiv 0 \bmod p^{n-a} ,
\]
whence
\[
\ord_p B + \sum_{\substack{\zeta^{p^{n-b}} = 1 \\ \zeta \ne 1}} \ord_p(A - \zeta) \ge n - a .
\]
But $A$ is $\equiv 1 \bmod 2p$, so for any non-trivial $p$-power root of unity we have
\[
\ord_p(A - \zeta) = \ord_p(1 - \zeta) .
\]
Therefore
\begin{align*}
\sum_{\substack{\zeta^{p^{n-b}} = 1 \\ \zeta \ne 1}} \ord_p(A - \zeta) &= \sum_{\substack{\zeta^{p^{n-b}} = 1 \\ \zeta \ne 1}} \ord_p(1 - \zeta) \\
&= \ord_p\left(\left.\left(\frac{X^{p^{n-b}}-1}{X-1}\right)\right|_{X=1}\right) \\
&= \ord_p(p^{n-b}) \\
&= n - b .
\end{align*}
This yields
\[
\ord_p(B) + n - b \ge n - a ,
\]
as required.
\end{proof}

\begin{krem}\label{rem:10.10.5}
Here is an example to show that the condition $(*)$ is essential. Consider the element
\[
\begin{pmatrix} -1 & 1 \\ 0 & 1 \end{pmatrix} \in \GL(2, \Z/4\Z) .
\]
It has order two, and the cyclic group $G$ it generates violates $(*)$. Consider the ring of invariants
\[
\Z[i, T]^G
\]
under the action
\[
\left\{
\begin{aligned}
i &\mapsto -i \\
T &\mapsto iT .
\end{aligned}
\right.
\]
An elementary calculation shows that the ring of invariants is
\begin{align*}
(\Z[i, T])^G &= \Z[T + iT, iT^2] \\
&= \Z[X, Y]/(X^2 = 2Y) ,
\end{align*}
which visibly is \emph{not regular} at the maximal ideal $(X = 0, Y = 0, 2 = 0)$.
\end{krem}

\begin{kcor}\label{cor:10.10.6}
Hypotheses and notations as in \eqref{thm:10.10.3}, let $k$ be a finite field of characteristic $p$, and $A \tens{\Z} B[1/N] \to k$ a ring homomorphism.
\par
Suppose that $\mathcal{P}[1/N]$ satisfies coarse base change near infinity (cf.~\ref{prop:8.6.6}, \ref{item:8.5.1}), (e.g., $\mathcal{P}[1/N]$ representable). Then the following conditions are equivalent.
\begin{enumerate}[label=(\arabic*)]
\item $\Gamma_1$ is balanced.
\item $\mathbf{M}(\mathcal{P}) \tens{A \tens{\Z} B} k$ is a smooth curve over $k$ outside the supersingular points [and the points where $j(j-1728)$ is not invertible].{\renewcommand{\thefootnote}{*}\footnote{The phrase in brackets can be deleted, see Notes Added in Proof.}}
\item $\mathbf{M}(\mathcal{P}) \tens{A \tens{\Z} B} k$ is generically reduced (i.e., reduced outside a finite set of points).
\end{enumerate}
\end{kcor}

\begin{proof}
By the theorem, $(1) \implies (2)$, and trivially $(2) \implies (3)$. Suppose that $\Gamma_1$ is not balanced. Then the calculation of the previous theorem shows that
\[
\mathcal{P}_{\Tate(q)/(A \tens{\Z} B[1/N])((q))}/\pm 1
\]
is a disjoint union of schemes, indexed by representatives $\Lambda$ of
\[
\HomSurj((\Z/p^n\Z)^2, \Z/p^n\Z)/\pm\Gamma_1 ,
\]
such that the $\Lambda$'th component is faithfully flat over $A_\Lambda \tens{\Z} B_\Lambda[1/N]$. By hypothesis, $\pm 1$ operates either freely or trivially on
\[
\mathcal{P}_{\Tate(q)/(A \tens{\Z} B[1/N])((q))} .
\]
Clearly we have
\[
((A \tens{\Z} B[1/N])((q))) \tens{A \tens{\Z} B[1/N]} k \xrightarrow{\sim} k((q)) ,
\]
so
\[
(\mathcal{P}_{\Tate(q)/k((q))})/\pm 1 = ((\mathcal{P}_{\Tate(q)/(A \tens{\Z} B[1/N])((q))})/\pm 1) \tens{} k
\]
is a disjoint union of schemes faithfully flat over the rings
\[
(A_\Lambda \tens{\Z} B_\Lambda[1/N]) \tens{A \tens{\Z} B[1/N]} k .
\]
If $\Gamma_1$ is not balanced at $\Lambda$, then $A_\Lambda$ is the ring of integers in a non-trivial extension of $A$ inside $\Z[\zeta_{p^n}]$, so it is fully ramified over $A$ at the unique place dividing $p$. Therefore
\[
(A_\Lambda \tens{\Z} B_\Lambda[1/N]) \tens{A \tens{\Z} B} k
\]
is \emph{not reduced}. Therefore, if $\Gamma_1$ is not balanced, the scheme
\[
\mathcal{P}_{\Tate(q)/k((q))}/\pm 1
\]
is not reduced. Because $\mathcal{P}[1/N]$ satisfies coarse base change, we have
\[
(\mathcal{P}_{\Tate(q)/k((q))})/\pm 1 \simeq \mathbf{M}(\mathcal{P} \tens{} k)_{k((q))} = (\mathbf{M}(\mathcal{P}) \tens{} k)_{k((q))} .
\]
Therefore $(\mathbf{M}(\mathcal{P}) \tens{} k)_{k((q))}$ is not reduced. But the field extension
\[
k(j) \hookrightarrow k((1/j)) = k((1/q))
\]
is separable, so $(\mathbf{M}(\mathcal{P}) \tens{} k)_{k(j)}$ is not reduced.
\par
Because $\mathbf{M}(\mathcal{P}) \tens{A \tens{\Z} B} k$ is finite over $\Spec(k[j])$, this means that $\mathbf{M}(\mathcal{P}) \tens{} k$ cannot be ``reduced outside a finite set of points.''
\end{proof}

\begin{kcor}[due to O.~Gabber]\label{cor:10.10.7}
Let $p_1, \dots, p_r$ be distinct primes, $k_1, \dots, k_r$ integers $\ge 1$, $N_0 \ge 1$ an integer prime to $p_1 \dots p_r$. Suppose we are given subgroups
\[
\left\{
\begin{aligned}
\Gamma_0 &\subset \SL(2, \Z/N_0\Z) \\
\Gamma_i &\subset \GL(2, \Z/p_i^{k_i}\Z) \quad\text{for } i = 1, \dots, r .
\end{aligned}
\right.
\]
Let $N = N_0 \prod p_i^{k_i}$, and denote by
\[
\Gamma \subset \GL(2, \Z/N\Z)
\]
the product subgroup
\[
\Gamma = \Gamma_0 \times \Gamma_1 \times \dots \times \Gamma_r .
\]
Suppose that both of the following conditions are satisfied.
\begin{enumerate}[label=(\alph*)]
\item For $i = 1, \dots, r$, the subgroup $\Gamma_i$ satisfies $(*)$ (cf.~\ref{eq:10.10.1.1}).
\item Either $-1$ operates without fixed points on the cusp-labels for $\Gamma$, or every one of the groups $\Gamma_0, \Gamma_1, \dots, \Gamma_r$ contains $-1$.
\end{enumerate}
Denote by $\mathcal{P}$ the moduli problem $[\Gamma(N)]/\Gamma$ on $(\mathrm{Ell}/\Z)$. Then the compactified coarse moduli scheme $\overline{\mathbf{M}}(\mathcal{P})$ is regular except possibly at points where $j(j-1728)$ is not invertible or at supersingular points in characteristic $p$ dividing $N$.
\end{kcor}

\begin{proof}
In the notation of \eqref{thm:10.9.7}, $\discrim(\Gamma)$ divides $p_1 \dots p_r$, and therefore $\overline{\mathbf{M}}(\mathcal{P})[1/p_1 \dots p_r]$ is smooth over the cyclotomic ring $(\Z[\zeta_N])^{\det(\Gamma)}[1/p_1 \dots p_r]$ outside the specified exceptional points. Therefore $\overline{\mathbf{M}}(\mathcal{P})[1/p_1 \dots p_r]$ is regular outside the specified exceptions. To see the regularity at points in characteristic $p_1, \dots, p_r$, we apply Theorem \ref{thm:10.10.3} to, say, $\Gamma_1$ and $\Gamma_0 \times \Gamma_2 \times \dots \times \Gamma_r$, which tells us that $\overline{\mathbf{M}}(\mathcal{P})[1/N_0 p_2 \dots p_r]$ is regular outside the specified exceptions.
\end{proof}

\begin{krem}\label{rem:10.10.8}
The example already cited (\ref{rem:10.10.5}) shows that the condition $(*)$ is essential. For if $p_1^{k_1} = 4$, and if we take for $\Gamma_1$ the subgroup of order two in $\GL(2, \Z/4\Z)$ generated by
\[
\begin{pmatrix} -1 & 1 \\ 0 & 1 \end{pmatrix} ,
\]
then $\overline{\mathbf{M}}(\mathcal{P})$ will be \emph{irregular} at every cusp in characteristic two.
\end{krem}

\section{Basic examples of balanced subgroups}
\label{sec:10.11}

\kmitem\label{item:10.11.1}
We again fix a prime number $p$, and an integer $n \ge 1$. For every subgroup $H$ of $(\Z/p^n\Z)^\times$, we denote by
\[
\begin{pmatrix} H & * \\ 0 & H \end{pmatrix} \subset \GL(2, \Z/p^n\Z)
\]
the subgroup consisting of all matrices of the form
\[
\begin{pmatrix} A & B \\ 0 & D \end{pmatrix} , \quad A, D \in H, \quad B \in \Z/p^n\Z .
\]

\begin{kprop}\label{prop:10.11.2}
The subgroup $\begin{pmatrix} H & * \\ 0 & H \end{pmatrix} \subset \GL(2, \Z/p^n\Z)$ is always partly balanced. It is balanced precisely for the following list of $H$:
\begin{enumerate}[label=(\arabic*)]
\item for $n = 1$, no restriction, i.e., $H = \text{any subgroup of } \mathbb{F}_p^\times$
\item for $n \ge 2$, and $p$ odd, or $n \ge 4$ and $p = 2$, $H$ must be one of the groups $1 + (p^a)$, for $a$ an integer $\ge [n/2]$
\item for $n = 2, 3$ and $p = 2$, $H$ is \emph{unrestricted}.
\end{enumerate}
\end{kprop}

\begin{proof}
Let $\Lambda \in \HomSurj((\Z/p^n\Z)^2, \Z/p^n\Z)$ correspond to a primitive vector $(x,y)$ (i.e., $\Lambda\begin{pmatrix} a \\ b \end{pmatrix} = ax + by$). We must compute, for each $\Lambda$, the group of determinants of
\[
\begin{pmatrix} H & * \\ 0 & H \end{pmatrix} \cap \Fix\Lambda .
\]
The condition that an element $\begin{pmatrix} A & B \\ 0 & D \end{pmatrix}$ fix $\Lambda$ is
\[
(x,y) = (x,y)\begin{pmatrix} A & B \\ 0 & D \end{pmatrix} = (Ax, Bx + Dy) ,
\]
i.e.,
\[
\left\{
\begin{aligned}
(A - 1)x &= 0 \\
(1 - D)y &= Bx .
\end{aligned}
\right.
\]
\par
\emph{Case 1.} $x$ invertible in $\Z/p^n\Z$. Then $A = 1$, but $D$ is unrestricted, because we may solve for $B = \left(\frac{1}{x}\right)(1-D)y$. Therefore any $\begin{pmatrix} H & * \\ 0 & H \end{pmatrix}$ is balanced at a component label $\Lambda$ with $x$ invertible.
\par
\emph{Case 2.} $x = 0$. Then $A$ is free, and we may take $D = 1$, $B$ arbitrary. So any $\begin{pmatrix} H & * \\ 0 & H \end{pmatrix}$ is balanced at a component label $\Lambda$ with $x = 0$.
\par
The first two cases complete the analysis for $n = 1$ (since $x \in \mathbb{F}_p$ is either invertible or zero).
\par
\emph{Case 3.} $n \ge 2$, $x = p^a \times (\text{unit})$ with $1 \le a \le n-1$. Then $y$ is invertible, because $(x,y)$ is a primitive vector. The conditions are
\[
\left\{
\begin{aligned}
A &\equiv 1 \pmod{p^{n-a}} \\
D &\equiv 1 \pmod{p^a} .
\end{aligned}
\right.
\]
Therefore for this $\Lambda = (x,y)$, we have
\[
\det\left(\begin{pmatrix} H & * \\ 0 & H \end{pmatrix} \cap (\Fix\Lambda)\right) = H \cap (1 + (p^{\min(a, n-a)})) .
\]
Therefore $\begin{pmatrix} H & * \\ 0 & H \end{pmatrix}$ is balanced at this $\Lambda$ if and only if
\[
H \subset 1 + (p^{\min(a, n-a)}) .
\]
The most restrictive choice of $a$ is $a = [n/2]$, i.e.,
\[
\begin{pmatrix} H & * \\ 0 & H \end{pmatrix} \text{ balanced} \iff H \equiv 1 \bmod p^{[n/2]} .
\]
Now for $p$ odd, the group $1 + (p) \subset (\Z/p^n\Z)^\times$ is cyclic of order $p^{n-1}$, and its only subgroups are the $1 + (p^\nu)$ with $1 \le \nu \le n$. Similarly, if $p = 2$ and $n \ge 4$, the group $1 + (4)$ is cyclic, and its only subgroups are the $1 + (2^{\nu+2})$'s, $0 \le \nu \le n-2$.
\end{proof}

\begin{krem}\label{rem:10.11.3}
For $p = 2$, the list of $H$ such that $\begin{pmatrix} H & * \\ 0 & H \end{pmatrix}$ is balanced \emph{and} satisfies $(*)$ is the following:
\begin{itemize}[label={}]
\item $n = 1$;\quad no restriction, any (!) subgroup of $\mathbb{F}_2^\times$
\item $n = 2$;\quad $H$ trivial
\item $n = 3$;\quad $H$ trivial or $H = 1 + (4) = \{1, 5\}$
\item $n \ge 4$;\quad $H = 1 + (2^a)$ with $a \ge [n/2]$.
\end{itemize}
\end{krem}

\section{Application to the moduli problems $[\Gamma_0(p^n); a, a]$}
\label{sec:10.12}

\kmitem\label{item:10.12.1}
We have already seen (\ref{thm:7.4.2}) that we have isomorphisms of moduli problems on $(\mathrm{Ell}/\Z)$
\[
[\Gamma(p^n)]/\begin{pmatrix} 1 & * \\ 0 & 1 \end{pmatrix} \simeq [\mathrm{bal.}\ \Gamma_1(p^n)] ,
\]
so that for any subgroup $H \subset (\Z/p^n\Z)^\times$, we have
\[
[\Gamma(p^n)]/\begin{pmatrix} H & * \\ 0 & H \end{pmatrix} \xrightarrow{\sim} [\mathrm{bal.}\ \Gamma_1(p^n)]/H \times H .
\]
We have seen (\ref{thm:7.6.1}) that this quotient is regular, finite and flat over $(\mathrm{Ell}/\Z)$. In the particular case
\[
H = 1 + (p^a) \subset (\Z/p^n\Z)^\times ,
\]
this quotient is (cf.~\ref{thm:7.9.6}) the moduli problem
\[
[\Gamma_0(p^n); a, a] .
\]

\begin{kthm}\label{thm:10.12.2}
Let $p$ be a prime, $n \ge 1$ an integer, $H \subset (\Z/p^n\Z)^\times$ one of the subgroups
\begin{enumerate}[label=(\arabic*)]
\item if $n = 1$, any subgroup of $\mathbb{F}_p^\times$
\item if $n \ge 2$, one of the subgroups $1 + (p^a)$ with $a \ge [n/2]$ for $p$ odd, and $a \ge \max(2, [n/2])$ if $p = 2$.
\end{enumerate}
Denote by $A$ the ring
\[
A = (\Z[\zeta_{p^n}])^H \subset \Z[\zeta_{p^n}] ,
\]
and by $\mathcal{P}$ the moduli problem
\[
([\mathrm{bal.}\ \Gamma_1(p^n)]/H \times H)^{A\text{-can}} \quad\text{on } (\mathrm{Ell}/A) .
\]
Let $N \ge 3$ be an integer prime to $p$, $\Gamma \subset \GL(2, \Z/N\Z)$ a subgroup, $B$ the ring
\[
B = \Z[\zeta_N]^{\det(\Gamma)} \subset \Z[\zeta_N] .
\]
Suppose that the moduli problem
\[
\mathcal{R} = ([\Gamma(N)]/\Gamma)^{B\text{-can}} \quad\text{on } (\mathrm{Ell}/B)
\]
has $\mathcal{R}[1/N]$ representable on $(\mathrm{Ell}/B[1/N])$, and that $-1$ operates without fixed points on the cusp-labels for $\Gamma$. Then
\begin{enumerate}[label=(\arabic*)]
\item The simultaneous problem $(\mathcal{P}, \mathcal{R}[1/N])$ on $(\mathrm{Ell}/A \tens{\Z} B[1/N])$ is representable.
\item Its compactified moduli scheme $\overline{\mathbf{M}}(\mathcal{P}, \mathcal{R}[1/N])$ is a regular two-dimensional scheme, finite and flat over $\overline{\mathbf{M}}(\mathcal{R}[1/N]) \tens{\Z} A$.
\item Outside the supersingular points in characteristic $p$,\linebreak
$\overline{\mathbf{M}}(\mathcal{P}, \mathcal{R}[1/N])$ is smooth of relative dimension one over $A \tens{\Z} B[1/N]$.
\item The scheme of cusps of $\overline{\mathbf{M}}(\mathcal{P}, \mathcal{R}[1/N])$ is finite etale over $A \tens{\Z} B[1/N]$.
\item Over $A \tens{\Z} B[1/Np]$, the scheme $\overline{\mathbf{M}}(\mathcal{P}, \mathcal{R}[1/N])[1/p]$ is a proper smooth curve with geometrically connected fibers.
\end{enumerate}
\end{kthm}

\begin{proof}
Apply the previous theorem (\ref{thm:10.10.3}) to the product group $\begin{pmatrix} H & * \\ 0 & H \end{pmatrix} \times \Gamma$, and remember that $H$ is chosen so that $\begin{pmatrix} H & * \\ 0 & H \end{pmatrix}$ is balanced. The point to remark is that the quotient of $[\Gamma(Np^n)]$ by the product group is indeed the simultaneous problem in question, as follows from (\ref{prop:7.3.1}) and (\ref{item:9.4.3.6}). The finite flatness of the projection map $\overline{\mathbf{M}}(\mathcal{P}, \mathcal{R}[1/N]) \to \overline{\mathbf{M}}(\mathcal{R}[1/N]) \tens{\Z} A$ results from the fact that it is a finite map (because it is a map of finite $\mathbb{P}^1_{A \tens{\Z} B[1/N]}$-schemes) between regular schemes of the same dimension.
\end{proof}

\begin{krem}\label{rem:10.12.3}
This theorem will be used later in studying questions of good reduction. In those applications, one very natural choice of auxiliary rigidifying moduli problem $\mathcal{R}$ will be $\Gamma_1(N)$, for an arbitrary integer $N \ge 4$ prime to $p$. In this case, the ring $B$ is simply $\Z$ itself. But $[\mathrm{bal.}\ \Gamma_1(N)]^{\mathrm{can}}$ with $N \ge 4$ prime to $p$, and its $B = \Z[\zeta_N]$, or $[\Gamma(N)]^{\mathrm{can}}$ with $N \ge 3$ prime to $p$ and $B = \Z[\zeta_N]$, are ``just as good'' for studying questions of good reduction at primes over $p$.
\end{krem}

\section{The numerology of modular schemes, via the line bundle \texorpdfstring{$\underline{\omega}$}{ω}}
\label{sec:10.13}

\kmitem\label{item:10.13.1}
Let $R$ be an excellent noetherian regular ring, $\mathcal{P}$ a representable moduli problem on $(\mathrm{Ell}/R)$ which is finite over $(\mathrm{Ell}/R)$, and normal near infinity. Over the moduli scheme $\mathfrak{M}(\mathcal{P})$ we have the corresponding universal elliptic curve
\[
\begin{tikzcd}[row sep=1.8em]
E \arrow[d, "f"] \\
\mathfrak{M}(\mathcal{P}) ,
\end{tikzcd}
\]
and the corresponding invertible sheaf $\underline{\omega}$ on $\mathfrak{M}(\mathcal{P})$ defined by
\begin{equation}\label{eq:10.13.1.1}
\underline{\omega} = f_*\Omega^1_{E/\mathfrak{M}(\mathcal{P})} = \Lie(E/\mathfrak{M}(\mathcal{P}))^\vee .
\end{equation}

\kmitem\label{item:10.13.2}
There is a canonical way to extend $\underline{\omega}^{\otimes 2}$ to an invertible sheaf, still denoted $\underline{\omega}^{\otimes 2}$, on the compactified moduli scheme $\overline{\mathfrak{M}}(\mathcal{P})$, as follows. It suffices to specify a prolongation of $\underline{\omega}^{\otimes 2}$ as invertible sheaf from
\[
\overline{\mathfrak{M}}(\mathcal{P})_{R((q))} \quad\text{to}\quad \overline{\mathfrak{M}}(\mathcal{P})_{R[[q]]} = \widehat{\Cusps}(\mathcal{P}) .
\]
We will do this by specifying a canonical \emph{trivialization}
\[
\mathcal{O} \xrightarrow{\ \sim\ } \underline{\omega}^{\otimes 2} \quad\text{restricted to}\quad \overline{\mathfrak{M}}(\mathcal{P})_{R((q))} ,
\]
and then patch together
\[
\left\{
\begin{aligned}
&\text{the given } \underline{\omega}^{\otimes 2} \text{ on } \mathfrak{M}(\mathcal{P}) \\
&\text{the trivial } \mathcal{O} \text{ on } \overline{\mathfrak{M}}(\mathcal{P})_{R[[q]]}
\end{aligned}
\right.
\]
along this trivialization. We construct the trivialization as follows. By (\ref{prop:8.11.7}), we have a cartesian diagram
\[
\begin{tikzcd}[row sep=2.0em, column sep=2.5em]
\Tate(q) \arrow[d] & \Tate(q) \arrow[l] \arrow[d] \arrow[r] & E_{\mathrm{univ}} \arrow[d] \\
\Spec(R((q))) & \mathcal{P}_{\Tate(q)/R((q))} \arrow[l, "f"'] \arrow[r, "\pi"] & \overline{\mathfrak{M}}(\mathcal{P})_{R((q))}
\end{tikzcd}
\]
in which $\pi$ is a finite etale $\pm 1$-torsor. Therefore $\pi^*(\underline{\omega})$ admits the nowhere vanishing section deduced by pull-back by $f$ of the canonical one-form (\ref{sec:8.8}, T.2)
\[
\omega_{\mathrm{can}} \quad\text{on}\quad \Tate(q)/\Z((q)) .
\]
Under the action of $\pm 1 = \Aut(\Tate(q)/R((q)))$, the one-form $\omega_{\mathrm{can}}$ is carried to $\pm \omega_{\mathrm{can}}$. Therefore its square $(\omega_{\mathrm{can}})^{\otimes 2}$ defines a trivialization of $\pi^*(\underline{\omega}^{\otimes 2})$ which is \emph{invariant} by $\pm 1$, so descends to give the sought-after trivialization of $\underline{\omega}^{\otimes 2}$ restricted to $\overline{\mathfrak{M}}(\mathcal{P})_{R((q))}$.

The same procedure serves to extend any even power $\underline{\omega}^{\otimes 2k}$ of $\underline{\omega}$ to an invertible sheaf on all of $\overline{\mathfrak{M}}(\mathcal{P})$. The extended $\underline{\omega}^{\otimes 2k}$ is related to the extended $\underline{\omega}^{\otimes 2}$ by the formula
\[
\underline{\omega}^{\otimes 2k}\text{ extended} = (\underline{\omega}^{\otimes 2}\text{ extended})^{\otimes k} ,
\]
simply because $\omega_{\mathrm{can}}^{\otimes 2k} = (\omega_{\mathrm{can}}^{\otimes 2})^{\otimes k}$.

\kmitem\label{item:10.13.3}
It may or may not be possible to extend $\underline{\omega}$ itself to a line bundle on all of $\overline{\mathfrak{M}}(\mathcal{P})$ whose square is the above extension of $\underline{\omega}^{\otimes 2}$. For an example where no such extension exists, consider the moduli problem $[\Delta = 1]$ on $(\mathrm{Ell}/\Z[1/6])$ defined by
\[
[\Delta = 1](E/R) = \left\{
\begin{aligned}
&\text{the set of nowhere-vanishing} \\
&\text{one-forms } \omega \text{ on } E/R \text{ such} \\
&\text{that } \Delta(E, \omega) = 1 .
\end{aligned}
\right.
\]
This is clearly a finite etale $\mu_{12}$-torsor over $(\mathrm{Ell}/\Z[1/6])$, which is represented by the smooth $\Z[1/6]$-curve in $(g_2, g_3)$-space
\[
(g_2)^3 - 27(g_3)^2 = 1 ,
\]
with universal $(E, \omega)$ given by
\[
(y^2 = 4x^3 - g_2 x - g_3,\ dx/y) .
\]
The scheme
\[
[\Delta = 1]_{\Tate(q)/\Z[1/6]((q))}
\]
is given by
\[
\Z[1/6][[q^{1/12}]] ,
\]
with universal object over it
\[
\Bigl(\Tate(q), \bigl(q^{1/12} \prod_{n \ge 1}(1 - q^n)^2\bigr) \omega_{\mathrm{can}}\Bigr) .
\]
The action of $\pm 1 = \Aut(\Tate(q))$ carries $\omega_{\mathrm{can}} \to \pm \omega_{\mathrm{can}}$, so it operates on $\Z[1/6][[q^{1/12}]]$ as $q^{1/12} \to \pm q^{1/12}$, with ring of invariants
\[
\widehat{\Cusps}([\Delta = 1]) = \Spec(\Z[1/6][[q^{1/6}]]) .
\]
The extended $\underline{\omega}^{\otimes 2}$ is free over $\widehat{\Cusps}([\Delta = 1])$ with nowhere vanishing section $\omega_{\mathrm{can}}^{\otimes 2}$. The square of the universal differential $\omega_{\mathrm{univ}}$ is equal to
\[
\Bigl(q^{1/6} \prod_{n \ge 1}(1 - q^n)^4\Bigr) (\omega_{\mathrm{can}})^{\otimes 2} ,
\]
which shows that $(\omega_{\mathrm{univ}})^{\otimes 2}$ extends to a global section of $\underline{\omega}^{\otimes 2}$ over $\overline{\mathfrak{M}}([\Delta = 1])$, which takes a simple zero at the cusp, and is elsewhere invertible. The scheme $\overline{\mathfrak{M}}([\Delta = 1])$ is just the complete elliptic curve over $\Z[1/6]$ whose affine equation is
\[
(g_2)^3 - 27(g_3)^2 = 1 ,
\]
and $\underline{\omega}^{\otimes 2}$ is the inverse ideal sheaf of the origin on this elliptic curve. In particular, it has degree one, so it is not the square of any line bundle!

In the positive direction, we have the following.

\begin{kprop}\label{prop:10.13.4}
Let $R$ be an excellent noetherian regular ring, $\mathcal{P}$ a representable moduli problem on $(\mathrm{Ell}/R)$ which is finite over $(\mathrm{Ell}/R)$ and normal near infinity. Let us denote by $Z(\mathcal{P})$ the normalization of $R[[q]]$ in the finite normal $R((q))$-scheme
\[
\mathcal{P}_{\Tate(q)/R((q))} .
\]
Suppose that the group $\pm 1 = \Aut(\Tate(q))$ operates freely on $Z(\mathcal{P})$. Then there is a natural way of extending $\underline{\omega}$ to a line bundle $\underline{\omega}$ on all of $\overline{\mathfrak{M}}(\mathcal{P})$, in such a way that
\[
(\underline{\omega}\text{ extended})^{\otimes 2} = \underline{\omega}^{\otimes 2}\text{ extended} .
\]
\end{kprop}

\begin{proof}
Because normalization commutes with quotients by finite groups, we have
\[
Z(\mathcal{P})/{\pm 1} \xrightarrow{\ \sim\ } \widehat{\Cusps}(\mathcal{P}) .
\]
Let $\mathcal{L}$ denote the trivial line bundle on $Z(\mathcal{P})$, which is the trivial extension of the trivial bundle
\[
\omega_{\mathrm{can}}\mathcal{O} \quad\text{on}\quad \mathcal{P}_{\Tate(q)/R((q))} .
\]
The group $\pm 1$ operates on $\mathcal{L}$, covering its action on $Z(\mathcal{P})$. Therefore if $\pm 1$ operates freely on $Z$, its action on $(Z(\mathcal{P}), \mathcal{L})$ defines a descent of $\mathcal{L}$ to $\widehat{\Cusps}(\mathcal{P})$, to a line bundle on $\widehat{\Cusps}(\mathcal{P})$ extending $\underline{\omega}$ from $\widehat{\Cusps}(\mathcal{P})_{R((q))}$. Because descent is functorial, the square of the descended bundle is the descent of the square $(\omega_{\mathrm{can}})^{\otimes 2}\mathcal{O}$, so this descent construction does indeed define a square root of $\underline{\omega}^{\otimes 2}$.
\end{proof}

\begin{kprop}\label{prop:10.13.5}
Let $R$ be an excellent noetherian regular ring, $\mathcal{P}$ and $\mathcal{P}'$ two representable moduli problems, both finite over $(\mathrm{Ell}/R)$ and both normal near infinity. Let
\[
f : \mathcal{P} \to \mathcal{P}'
\]
be a morphism of moduli problems over $(\mathrm{Ell}/R)$. Then
\begin{enumerate}[label=(\arabic*)]
\item Under the induced map of modular schemes $f : \mathfrak{M}(\mathcal{P}) \to \mathfrak{M}(\mathcal{P}')$ we have
\[
f^*(\underline{\omega}) \simeq \underline{\omega} .
\]
\item Under the induced map of compactified modular schemes we have
\[
\bar{f} : \overline{\mathfrak{M}}(\mathcal{P}) \to \overline{\mathfrak{M}}(\mathcal{P}') ,
\]
\[
\bar{f}^*(\underline{\omega}^{\otimes 2}) = \underline{\omega}^{\otimes 2} .
\]
\item If $\pm 1$ acts freely on $Z(\mathcal{P}')$, then it operates freely on $Z(\mathcal{P})$, and under the induced map of compactified modular schemes we have
\[
\bar{f}^*(\underline{\omega}) = \underline{\omega} .
\]
\end{enumerate}
\end{kprop}

\begin{proof}
Assertion (1) is simply the fact that formation of $\underline{\omega}_{E/S}$ commutes with arbitrary change of base $S' \to S$. Assertion (2) is clear from the construction of the extension of $\underline{\omega}^{\otimes 2}$ via the differential $\omega_{\mathrm{can}}$ on $\Tate(q)$. Assertion (3) holds because $f$ induces a map $Z(\mathcal{P}) \to Z(\mathcal{P}')$ which is $\pm 1$ equivariant, so $\pm 1$ operates ``more'' freely on $Z(\mathcal{P})$ than on $Z(\mathcal{P}')$, and by the functoriality of descent.
\end{proof}

\begin{kprop}\label{prop:10.13.6}
Hypotheses and notations as in (\ref{item:10.13.1}) above, suppose that $\overline{\mathfrak{M}}(\mathcal{P})$ is smooth over $R$ in a neighborhood of the cusps, and that the cusps are finite etale over $R$. Then for any extension of scalars $R \to R'$ of excellent noetherian regular rings, the isomorphism (\ref{prop:8.6.6})
\[
\overline{\mathfrak{M}}(\mathcal{P}) \tens{R} R' \xrightarrow{\ \sim\ } \overline{\mathfrak{M}}(\mathcal{P} \tens{R} R')
\]
induces an isomorphism of line bundles
\[
\underline{\omega}^{\otimes 2}_{\overline{\mathfrak{M}}(\mathcal{P})} \tens{R} R' \xrightarrow{\ \sim\ } \underline{\omega}^{\otimes 2}_{\overline{\mathfrak{M}}(\mathcal{P} \tens{R} R')} ,
\]
and, if $\pm 1$ operates freely on $Z(\mathcal{P})$, an isomorphism
\[
\underline{\omega}_{\overline{\mathfrak{M}}(\mathcal{P})} \tens{R} R' \xrightarrow{\ \sim\ } \underline{\omega}_{\overline{\mathfrak{M}}(\mathcal{P} \tens{R} R')} .
\]
\end{kprop}

\begin{proof}
Again clear from the definitions.
\end{proof}

\begin{kprop}\label{prop:10.13.7}
Let $N \ge 1$ be an integer, $\Gamma \subset \GL(2, \allowbreak\Z/N\Z)$ a subgroup, $A$ the ring $(\Z[\zeta_N])^{\det(\Gamma)}$, $B$ an excellent noetherian regular $A$-algebra; $\mathcal{P}$ the moduli problem
\[
([\Gamma(N)]/\Gamma)^{A\text{\rm-can}} \tens{A} B \quad\text{on } (\mathrm{Ell}/B) .
\]
Suppose that $\mathcal{P}$ is normal near infinity, so that the scheme $Z(\mathcal{P}) =$ the normalization of $B[[q]]$ in
\[
\mathcal{P}_{\Tate(q)/B((q))}
\]
makes sense. If $-1$ operates without fixed points on the space
\[
\HomSurj((\Z/N\Z)^2, \Z/N\Z)/\Gamma
\]
of cusp-labels for $\Gamma$, then $\pm 1$ operates freely on $Z(\mathcal{P})$.
\end{kprop}

\begin{proof}
By (\ref{thm:10.8.2}), the scheme $\mathcal{P}_{\Tate(q)/B((q))}$ is a disjoint union of schemes, indexed by $\HomSurj/\Gamma$, and $\pm 1$ operates compatibly with its action on $\HomSurj/\Gamma$. Therefore $Z(\mathcal{P})$ is a disjoint union over the same indexing set, $\HomSurj/\Gamma$, with $\pm 1$ acting on $Z(\mathcal{P})$ compatibly with its action on $\HomSurj/\Gamma$. So if $\pm 1$ operates freely on $\HomSurj/\Gamma$, it certainly operates freely on $Z(\mathcal{P})$.
\end{proof}

\begin{kcor}\label{cor:10.13.8}
Hypotheses and notations as above, suppose that $B$ is an $A[1/N]$-algebra. Denote by $\Gamma_0$ the intersection $\Gamma \cap \SL(2, \Z/N\Z)$. If $-1$ operates without fixed points on the space $\HomSurj/\Gamma_0$ of cusp-labels for $\Gamma_0$, then $\pm 1$ operates freely on $Z(\mathcal{P})$.
\end{kcor}

\begin{proof}
It suffices to check that $\pm 1$ acts freely on $Z(\mathcal{P})$ after the finite etale base-change $B \mapsto B' = B \tens{A[1/N]} \Z[\zeta_N, 1/N]$, under which
\[
Z(\mathcal{P}) \tens{B} B' \xrightarrow{\ \sim\ } Z(\mathcal{P} \tens{B} B') = Z(\mathcal{P}')
\]
where $\mathcal{P}'$ is the moduli problem attached to $\Gamma_0 = \Gamma \cap \SL(2, \Z/N\Z)$
\[
\mathcal{P}' = ([\Gamma(N)]/\Gamma_0)^{\mathrm{can}} \tens{\Z[\zeta_N]} B' \quad\text{on } (\mathrm{Ell}/B') .
\]
Now apply the previous proposition to $\mathcal{P}'$.
\end{proof}

\begin{kcor}\label{cor:10.13.9}
Hypotheses and notations as above, suppose that $B$ is an $A[1/N]$-algebra, and that $\mathcal{P}$ is representable. If $N$ is odd, or if $N = 1, 2, 4$ and $\Gamma_0 = \Gamma \cap \SL(2, \Z/N\Z)$ lies in the upper unipotent subgroup $\begin{pmatrix} 1 & * \\ 0 & 1 \end{pmatrix}$, then $\pm 1$ operates freely on $Z(\mathcal{P})$, and consequently $\underline{\omega}$ ``exists'' on $\overline{\mathfrak{M}}(\mathcal{P})$.
\end{kcor}

\begin{proof}
Simply apply (\ref{lem:10.7.2}) or (\ref{lem:10.7.4}) to $\Gamma_0$.
\end{proof}

\par\smallskip\noindent
\textbf{(10.13.9.1)}\quad\emph{Summarizing Table}
\begin{center}
\textsc{Short Table of Applications}
\vspace{0.5em}

{\setlength{\tabcolsep}{4.5pt}%
\begin{tabular}{c|c|c}
\hline
$\mathcal{P}$ & representable over & $\underline{\omega}$ exists on $\overline{\mathfrak{M}}(\mathcal{P})$ \\
\hline
$\Gamma(N)^{\mathrm{can}}$, $N \ge 3$ & any $\Z[\zeta_N, 1/N]$-algebra & yes \\
\hline
$\Gamma_1(N)^{\mathrm{can}}$, $N \ge 5$ & any $\Z[\zeta_N, 1/N]$-algebra & yes \\
\hline
$\Gamma_1(N)$, $N \ge 5$ & any $\Z[1/N]$-algebra & yes \\
\hline
\end{tabular}}
\end{center}

\kmitem\label{item:10.13.10}
Recall that if $E/S/R$ is an elliptic curve over a smooth $R$-scheme $S$, $E \xrightarrow{f} S$, the Kodaira-Spencer mapping (cf.~\cite[A1.4]{K2}) is an $\mathcal{O}_S$-morphism
\[
\begin{aligned}
T_{S/R} &\to \mathrm{R}^1 f_*(T_{E/S}) \simeq \mathrm{R}^1 f_*(\mathcal{O}) \tens{} \Lie(E/S) \\
&\simeq \Lie(E/S) \tens{} \Lie(E/S) \\
&\simeq (\underline{\omega}^2_{E/S})^\vee ,
\end{aligned}
\]
whose $\mathcal{O}_S$-linear dual is an $\mathcal{O}_S$ linear map
\[
(\underline{\omega}_{E/S})^{\otimes 2} \to \Omega^1_{S/R} .
\]
One knows that this mapping is an isomorphism if and only if the representable moduli problem on $(\mathrm{Ell}/R)$ represented by $E/S/R$ is etale over $(\mathrm{Ell}/R)$.

Thus if $\mathcal{P}$ is a representable problem which is etale over $(\mathrm{Ell}/R)$, we have a canonical isomorphism of line bundles on $\mathfrak{M}(\mathcal{P})$
\[
\underline{\omega}^{\otimes 2} \xrightarrow{\ \sim\ } \Omega^1_{\mathfrak{M}(\mathcal{P})/R} .
\]

\begin{kthm}\label{thm:10.13.11}
Let $R$ be an excellent noetherian regular domain whose fraction field has characteristic zero. Let $\mathcal{P}$ be a representable moduli problem on $(\mathrm{Ell}/R)$ which is finite etale over $(\mathrm{Ell}/R)$. [Then by (\ref{thm:8.6.8}), we know that $\overline{\mathfrak{M}}(\mathcal{P})$ is a proper smooth curve over $R$, in which the cusps are finite etale over $R$ and for which $d\log(1/j)$ is a basis, near the cusps, for $\Omega^1_{\overline{\mathfrak{M}}(\mathcal{P})/R}(\log\mathrm{cusps})$.] The Kodaira-Spencer isomorphism on $\mathfrak{M}(\mathcal{P})$
\[
\underline{\omega}^{\otimes 2} \xrightarrow{\ \sim\ } \Omega^1_{\mathfrak{M}(\mathcal{P})/R}
\]
extends to an isomorphism on $\overline{\mathfrak{M}}(\mathcal{P})$
\[
\underline{\omega}^{\otimes 2} \xrightarrow{\ \sim\ } \Omega^1_{\overline{\mathfrak{M}}(\mathcal{P})/R}(\log\mathrm{cusps}) ,
\]
where $\underline{\omega}^{\otimes 2}$ denotes the canonical extension described above of $\underline{\omega}^{\otimes 2}$ on $\mathfrak{M}(\mathcal{P})$. For any regular noetherian excellent $R$-algebra $R'$, this remains true after the base-change $R \to R'$.
\end{kthm}

\begin{proof}
In view of (\ref{thm:8.6.8}), and the formula
\[
1/j = q \times (\text{invertible series in } \Z[[q]]) ,
\]
we see that in $\widehat{\Cusps}(\mathcal{P})$, $\Omega^1(\log\mathrm{cusps})$ admits as basis $d\log(q)$. In view of the definition of the extension $\underline{\omega}^{\otimes 2}$ in terms of the square $(\omega_{\mathrm{can}})^{\otimes 2}$ of the canonical differential $\omega_{\mathrm{can}}$ on $\Tate(q)$, it suffices to remark that by an explicit computation (cf.~\cite{K2}), one knows that for the Tate curve $\Tate(q)/\Z((q))$, the Kodaira-Spencer isomorphism carries
\[
(\omega_{\mathrm{can}})^{\otimes 2} \mapsto \frac{dq}{q} . \qedhere
\]
\end{proof}

\noindent\textbf{Remark.}
We will see later how ``wrong'' this theorem can become in the presence of ramification at finite distance, when the $\mathcal{P}$ in question does not ``come from characteristic zero'' (cf.~\ref{item:12.9.1}).

\begin{kcor}\label{cor:10.13.12}
Let $N \ge 1$ be an integer, $\Gamma \subset \GL(2, \allowbreak\Z/N\Z)$ a subgroup, $A$ the ring $(\Z[\zeta_N])^{\det\Gamma}$, $\mathcal{P}$ the moduli problem
\[
([\Gamma(N)]/\Gamma)^{A\text{\rm-can}} \tens{A} A[1/N] \quad\text{on } (\mathrm{Ell}/A[1/N]) .
\]
Suppose $\mathcal{P}$ is representable. Then $\overline{\mathfrak{M}}(\mathcal{P})$ is a proper smooth curve over $A[1/N]$ with geometrically connected fibers, in which the cusps are finite etale over $A[1/N]$. Let us denote by
\begin{itemize}[leftmargin=2.5em, itemsep=2pt, topsep=3pt]
\item[] $\Gamma_0 = \Gamma \cap \SL(2, \Z/N\Z)$
\item[] $g(\mathcal{P}) = {}$the common genus of the geometric fibers of $\overline{\mathfrak{M}}(\mathcal{P})$ over $A[1/N]$
\item[] $c(\mathcal{P}) = {}$the (common) number of cusps on the geometric fibers of $\overline{\mathfrak{M}}(\mathcal{P})$ over $A[1/N]$
\item[] $\deg(\mathcal{P}) = {}$the degree with which $\mathcal{P}$ is finite etale over $(\mathrm{Ell})$
\item[] \phantom{$\deg(\mathcal{P})$} ${}= 2 \times {}$(the degree of $\mathfrak{M}(\mathcal{P})$ as a finite flat covering of $\Spec(A[1/N][j]))$
\item[] $\deg(\underline{\omega}_{\mathcal{P}}) = {}$one-half the common degree of the restriction to any geometric fiber of the line bundle $\underline{\omega}^{\otimes 2}$ on $\overline{\mathfrak{M}}(\mathcal{P})$ over $A[1/N]$
\item[] \phantom{$\deg(\underline{\omega}_{\mathcal{P}})$} ${}= {}$the degree of $\underline{\omega}$ when that exists on $\overline{\mathfrak{M}}(\mathcal{P})$.
\end{itemize}
These quantities are related by the following formulas:
\[
\left\{
\begin{aligned}
2\deg(\underline{\omega}_{\mathcal{P}}) &= 2g(\mathcal{P}) - 2 + c(\mathcal{P}) \\
c(\mathcal{P}) &= \#(\HomSurj((\Z/N\Z)^2, \Z/N\Z)/\pm \Gamma_0) \\
\deg(\mathcal{P}) &= \#(\SL(2, \Z/N\Z)/\Gamma_0) \\
\deg(\underline{\omega}_{\mathcal{P}}) &= \frac{1}{24}\deg(\mathcal{P}) = \frac{\#(\SL(2, \Z/N\Z))}{24\#(\Gamma_0)} .
\end{aligned}
\right.
\]
\end{kcor}

\begin{proof}
Extending scalars from $A[1/N]$ to $\Z[\zeta_N, 1/N]$ amounts to replacing $\Gamma$ by $\Gamma_0 = \Gamma \cap \SL(2, \Z/N\Z)$. The first formula is the numerical incarnation of the isomorphism
\[
\underline{\omega}^{\otimes 2} \simeq \Omega^1(\log\mathrm{cusps}) \quad\text{on } \overline{\mathfrak{M}}(\mathcal{P}) .
\]
The second has already been proven (\ref{thm:10.9.1}). The third is clear since $[\Gamma(N)]^{\mathrm{can}}[1/N]$ is finite etale galois over $(\mathrm{Ell}/\Z[\zeta_N, 1/N])$ with galois group $\SL(2, \Z/N\Z)$, and $\mathcal{P}$ is its quotient by $\Gamma_0$. To prove the fourth, we argue as follows. By explicit calculation (\cite{Ig2} and \cite[AI, 5.5]{K2}) the theorem is true for $N = 3$, and $\Gamma_0 = \{1\}$; in this case
\[
\left\{
\begin{aligned}
&\overline{\mathfrak{M}} = \mathbb{P}^1, \text{ the } \mu\text{-line} \\
&\text{cusps} = \{\infty\} \cup \{\text{the roots of } \mu^3 = 27\} \\
&\underline{\omega} = \mathcal{O}(1) \\
&\text{universal curve: } X^3 + Y^3 + Z^3 = \mu XYZ .
\end{aligned}
\right.
\]
By (\ref{prop:10.13.5}), the ratio
\[
\frac{\deg(\underline{\omega}_{\mathcal{P}})}{\deg(\mathcal{P})}
\]
is independent of the particular choice of $\mathcal{P}$ of the form $[\Gamma(N)]^{\mathrm{can}}/\Gamma_0$ for some $N$ and some $\Gamma_0 \subset \SL(2, \Z/N\Z)$, for (\ref{prop:10.13.5}) allows us to replace $\mathcal{P}$ by $[\Gamma(N)]^{\mathrm{can}}$, $[\Gamma(N)]^{\mathrm{can}}$ by $[\Gamma(3N)]^{\mathrm{can}}$, and finally $[\Gamma(3)]^{\mathrm{can}}$, in which case we have checked it directly.
\end{proof}
